% Tactic found: the paper. The prose follows docs/writing.md. Every number,
% table, figure and board is written by `lab paper tactic-found` into
% generated/, from the model's committed results, and read back with \V{key}
% and \block{name}.
\documentclass[11pt]{article}
% The preamble every paper of the lab shares: % The preamble every paper of the lab shares: % The preamble every paper of the lab shares: \input{../../../papers/preamble}
% from models/<model>/paper/, \input{../preamble} from papers/<name>/.
% Boards come from chessboard and xskak, plots from pgfplots. A paper prints no
% number by hand: `lab paper` writes every number, table and figure under the
% paper's generated/ folder, and paper.tex says what they mean.
\usepackage[LSBC4,T1]{fontenc}
\usepackage[utf8]{inputenc}
\usepackage{lmodern}
\usepackage[margin=1in]{geometry}
\usepackage{microtype}
\usepackage{amsmath}
\usepackage{booktabs}
\usepackage{tabularx}
\usepackage{array}
\usepackage{xcolor}
\usepackage[obeyspaces,spaces]{url}
\usepackage{xskak}
\usepackage{pgfplots}
\usepgfplotslibrary{groupplots}
\usetikzlibrary{positioning,arrows.meta,calc,fit}
\pgfplotsset{compat=1.18}
\usepackage{caption}
\usepackage{subcaption}
\usepackage{enumitem}
\usepackage[section]{placeins}
\usepackage[hidelinks]{hyperref}

% A generated value by its key: generated/values.tex sets them, \V{key} prints
% one, and a key nobody set stops the build.
\makeatletter
\newcommand{\setvalue}[2]{\expandafter\gdef\csname labvalue@#1\endcsname{#2}}
\newcommand{\V}[1]{\@ifundefined{labvalue@#1}%
  {\PackageError{lab}{No generated value named '#1'}{Run `lab paper` for this model again.}}%
  {\csname labvalue@#1\endcsname}}
\makeatother

% Chart colours. A run keeps its colour in every figure of every paper.
\definecolor{labBlue}{HTML}{2A78D6}
\definecolor{labOrange}{HTML}{EB6834}
\definecolor{labAqua}{HTML}{1BAF7A}
\definecolor{labYellow}{HTML}{EDA100}
\definecolor{labMagenta}{HTML}{E87BA4}
\definecolor{labGreen}{HTML}{008300}
\definecolor{labViolet}{HTML}{6250D6}
\definecolor{labGray}{HTML}{898781}
\definecolor{labInk}{HTML}{0B0B0B}
\definecolor{labSecondary}{HTML}{52514E}
\definecolor{labGrid}{HTML}{E1E0D9}
\definecolor{labAxis}{HTML}{C3C2B7}
\definecolor{labWash}{HTML}{CDE2FB}

\pgfplotsset{
  lab/.style={
    font=\sffamily\footnotesize,
    axis line style={labAxis},
    tick style={labAxis},
    tick label style={labSecondary},
    label style={labSecondary},
    title style={labInk, font=\sffamily\small\bfseries, align=left},
    grid style={labGrid, line width=0.3pt},
    major grid style={labGrid, line width=0.3pt},
    axis lines*=left,
    tick align=outside,
    clip=false,
    legend style={draw=none, fill=none, font=\sffamily\footnotesize, text=labSecondary,
      cells={anchor=west}, legend columns=-1, column sep=6pt, at={(0,1.03)}, anchor=south west},
    every axis plot/.append style={line width=1.1pt, line join=round, line cap=round},
    scaled ticks=false,
    /pgf/number format/1000 sep={,},
  },
  lab marks/.style={mark=*, mark size=1.9pt, mark options={draw=white, line width=0.6pt}},
}

\setchessboard{boardfontencoding=LSBC4, boardfontfamily=skaknew, blackfieldcolor=black!18,
  whitefieldcolor=white, setfontcolors, showmover=false, label=true, labelfontsize=6pt,
  boardfontsize=12pt}

% A move or line in UCI, a position in FEN (breakable at a slash), and a run's name.
\newcommand{\uci}[1]{\texttt{#1}}
\newcommand{\fen}[1]{{\footnotesize\path{#1}}}
\newcommand{\run}[1]{\texttt{#1}}
\newcolumntype{L}{>{\raggedright\arraybackslash}X}
\newcolumntype{R}{>{\raggedleft\arraybackslash}X}

\setlength{\parskip}{0.4em}
\setlength{\parindent}{0pt}
\captionsetup{font=small, labelfont=bf}

% from models/<model>/paper/, \input{../preamble} from papers/<name>/.
% Boards come from chessboard and xskak, plots from pgfplots. A paper prints no
% number by hand: `lab paper` writes every number, table and figure under the
% paper's generated/ folder, and paper.tex says what they mean.
\usepackage[LSBC4,T1]{fontenc}
\usepackage[utf8]{inputenc}
\usepackage{lmodern}
\usepackage[margin=1in]{geometry}
\usepackage{microtype}
\usepackage{amsmath}
\usepackage{booktabs}
\usepackage{tabularx}
\usepackage{array}
\usepackage{xcolor}
\usepackage[obeyspaces,spaces]{url}
\usepackage{xskak}
\usepackage{pgfplots}
\usepgfplotslibrary{groupplots}
\usetikzlibrary{positioning,arrows.meta,calc,fit}
\pgfplotsset{compat=1.18}
\usepackage{caption}
\usepackage{subcaption}
\usepackage{enumitem}
\usepackage[section]{placeins}
\usepackage[hidelinks]{hyperref}

% A generated value by its key: generated/values.tex sets them, \V{key} prints
% one, and a key nobody set stops the build.
\makeatletter
\newcommand{\setvalue}[2]{\expandafter\gdef\csname labvalue@#1\endcsname{#2}}
\newcommand{\V}[1]{\@ifundefined{labvalue@#1}%
  {\PackageError{lab}{No generated value named '#1'}{Run `lab paper` for this model again.}}%
  {\csname labvalue@#1\endcsname}}
\makeatother

% Chart colours. A run keeps its colour in every figure of every paper.
\definecolor{labBlue}{HTML}{2A78D6}
\definecolor{labOrange}{HTML}{EB6834}
\definecolor{labAqua}{HTML}{1BAF7A}
\definecolor{labYellow}{HTML}{EDA100}
\definecolor{labMagenta}{HTML}{E87BA4}
\definecolor{labGreen}{HTML}{008300}
\definecolor{labViolet}{HTML}{6250D6}
\definecolor{labGray}{HTML}{898781}
\definecolor{labInk}{HTML}{0B0B0B}
\definecolor{labSecondary}{HTML}{52514E}
\definecolor{labGrid}{HTML}{E1E0D9}
\definecolor{labAxis}{HTML}{C3C2B7}
\definecolor{labWash}{HTML}{CDE2FB}

\pgfplotsset{
  lab/.style={
    font=\sffamily\footnotesize,
    axis line style={labAxis},
    tick style={labAxis},
    tick label style={labSecondary},
    label style={labSecondary},
    title style={labInk, font=\sffamily\small\bfseries, align=left},
    grid style={labGrid, line width=0.3pt},
    major grid style={labGrid, line width=0.3pt},
    axis lines*=left,
    tick align=outside,
    clip=false,
    legend style={draw=none, fill=none, font=\sffamily\footnotesize, text=labSecondary,
      cells={anchor=west}, legend columns=-1, column sep=6pt, at={(0,1.03)}, anchor=south west},
    every axis plot/.append style={line width=1.1pt, line join=round, line cap=round},
    scaled ticks=false,
    /pgf/number format/1000 sep={,},
  },
  lab marks/.style={mark=*, mark size=1.9pt, mark options={draw=white, line width=0.6pt}},
}

\setchessboard{boardfontencoding=LSBC4, boardfontfamily=skaknew, blackfieldcolor=black!18,
  whitefieldcolor=white, setfontcolors, showmover=false, label=true, labelfontsize=6pt,
  boardfontsize=12pt}

% A move or line in UCI, a position in FEN (breakable at a slash), and a run's name.
\newcommand{\uci}[1]{\texttt{#1}}
\newcommand{\fen}[1]{{\footnotesize\path{#1}}}
\newcommand{\run}[1]{\texttt{#1}}
\newcolumntype{L}{>{\raggedright\arraybackslash}X}
\newcolumntype{R}{>{\raggedleft\arraybackslash}X}

\setlength{\parskip}{0.4em}
\setlength{\parindent}{0pt}
\captionsetup{font=small, labelfont=bf}

% from models/<model>/paper/, \input{../preamble} from papers/<name>/.
% Boards come from chessboard and xskak, plots from pgfplots. A paper prints no
% number by hand: `lab paper` writes every number, table and figure under the
% paper's generated/ folder, and paper.tex says what they mean.
\usepackage[LSBC4,T1]{fontenc}
\usepackage[utf8]{inputenc}
\usepackage{lmodern}
\usepackage[margin=1in]{geometry}
\usepackage{microtype}
\usepackage{amsmath}
\usepackage{booktabs}
\usepackage{tabularx}
\usepackage{array}
\usepackage{xcolor}
\usepackage[obeyspaces,spaces]{url}
\usepackage{xskak}
\usepackage{pgfplots}
\usepgfplotslibrary{groupplots}
\usetikzlibrary{positioning,arrows.meta,calc,fit}
\pgfplotsset{compat=1.18}
\usepackage{caption}
\usepackage{subcaption}
\usepackage{enumitem}
\usepackage[section]{placeins}
\usepackage[hidelinks]{hyperref}

% A generated value by its key: generated/values.tex sets them, \V{key} prints
% one, and a key nobody set stops the build.
\makeatletter
\newcommand{\setvalue}[2]{\expandafter\gdef\csname labvalue@#1\endcsname{#2}}
\newcommand{\V}[1]{\@ifundefined{labvalue@#1}%
  {\PackageError{lab}{No generated value named '#1'}{Run `lab paper` for this model again.}}%
  {\csname labvalue@#1\endcsname}}
\makeatother

% Chart colours. A run keeps its colour in every figure of every paper.
\definecolor{labBlue}{HTML}{2A78D6}
\definecolor{labOrange}{HTML}{EB6834}
\definecolor{labAqua}{HTML}{1BAF7A}
\definecolor{labYellow}{HTML}{EDA100}
\definecolor{labMagenta}{HTML}{E87BA4}
\definecolor{labGreen}{HTML}{008300}
\definecolor{labViolet}{HTML}{6250D6}
\definecolor{labGray}{HTML}{898781}
\definecolor{labInk}{HTML}{0B0B0B}
\definecolor{labSecondary}{HTML}{52514E}
\definecolor{labGrid}{HTML}{E1E0D9}
\definecolor{labAxis}{HTML}{C3C2B7}
\definecolor{labWash}{HTML}{CDE2FB}

\pgfplotsset{
  lab/.style={
    font=\sffamily\footnotesize,
    axis line style={labAxis},
    tick style={labAxis},
    tick label style={labSecondary},
    label style={labSecondary},
    title style={labInk, font=\sffamily\small\bfseries, align=left},
    grid style={labGrid, line width=0.3pt},
    major grid style={labGrid, line width=0.3pt},
    axis lines*=left,
    tick align=outside,
    clip=false,
    legend style={draw=none, fill=none, font=\sffamily\footnotesize, text=labSecondary,
      cells={anchor=west}, legend columns=-1, column sep=6pt, at={(0,1.03)}, anchor=south west},
    every axis plot/.append style={line width=1.1pt, line join=round, line cap=round},
    scaled ticks=false,
    /pgf/number format/1000 sep={,},
  },
  lab marks/.style={mark=*, mark size=1.9pt, mark options={draw=white, line width=0.6pt}},
}

\setchessboard{boardfontencoding=LSBC4, boardfontfamily=skaknew, blackfieldcolor=black!18,
  whitefieldcolor=white, setfontcolors, showmover=false, label=true, labelfontsize=6pt,
  boardfontsize=12pt}

% A move or line in UCI, a position in FEN (breakable at a slash), and a run's name.
\newcommand{\uci}[1]{\texttt{#1}}
\newcommand{\fen}[1]{{\footnotesize\path{#1}}}
\newcommand{\run}[1]{\texttt{#1}}
\newcolumntype{L}{>{\raggedright\arraybackslash}X}
\newcolumntype{R}{>{\raggedleft\arraybackslash}X}

\setlength{\parskip}{0.4em}
\setlength{\parindent}{0pt}
\captionsetup{font=small, labelfont=bf}

% Written by `lab paper tactic-found`; do not edit by hand.
\looktrue
\chancetrue
\sequencefalse
\contextgbdttrue
\pairstrue
\variantstrue
\shapestrue
\gamestrue
\featurestrue
\changestrue
\changehandtrue
\changeepochstrue
\changestarttrue
\changeattackstrue
\changegeometrytrue
\changeconvtrue
\changehybridtrue
\changelargerfalse
\moretrue
\shownfalse
\shownpairedfalse
\setvalue{band.1.floor}{0.6428}
\setvalue{band.1.name}{\textless{}1700}
\setvalue{band.4.floor}{0.5989}
\setvalue{band.4.name}{2100-2299}
\setvalue{bands.best-ahead.of}{all five bands}
\setvalue{bands.count}{five}
\setvalue{bands.lead.narrowest}{+0.00852}
\setvalue{bands.lead.narrowest.name}{2100-2299}
\setvalue{bands.lead.widest}{+0.0162}
\setvalue{bands.lead.widest.name}{\textless{}1700}
\setvalue{best.beside.closer}{56}
\setvalue{best.beside.correlation}{0.83}
\setvalue{best.bottom.1.rate}{86.3}
\setvalue{best.bottom.10.lift}{2.2}
\setvalue{best.bottom.10.rate}{69.2}
\setvalue{best.bottom.10.recall}{22.2}
\setvalue{best.brier}{0.1778}
\setvalue{best.calibration-error}{0.0047}
\setvalue{best.calibration.groups}{ten}
\setvalue{best.calibration.highest}{96.1}
\setvalue{best.calibration.lowest}{30.2}
\setvalue{best.calibration.widest-gap}{1.3}
\setvalue{best.device}{NVIDIA A100-SXM4-40GB}
\setvalue{best.epochs.run}{11}
\setvalue{best.gain.of-lift}{16}
\setvalue{best.gain.percent}{2.3}
\setvalue{best.name}{\run{transformer-medium-hand-pretrained}}
\setvalue{best.parameters}{6{,}565{,}953}
\setvalue{best.platt.slope}{0.92}
\setvalue{best.pr-auc}{0.8667}
\setvalue{best.roc-auc}{0.7513}
\setvalue{best.rows-a-second}{14{,}000}
\setvalue{best.score}{0.5287}
\setvalue{best.score.high}{0.5321}
\setvalue{best.score.low}{0.5254}
\setvalue{best.seconds}{9 min}
\setvalue{best.top.1.rate}{99.8}
\setvalue{best.top.10.lift}{1.4}
\setvalue{best.top.10.rate}{95.9}
\setvalue{challenger.name}{\run{transformer-medium-hand-pretrained}}
\setvalue{code.commit}{\texttt{36a0ad9}}
\setvalue{context.lift.share}{35}
\setvalue{cost.seconds}{76 min}
\setvalue{cost.seconds.networks}{73 min}
\setvalue{cost.seconds.others}{3 min}
\setvalue{data.fingerprint}{\texttt{15bd888b36b8}}
\setvalue{data.licence}{CC0 1.0}
\setvalue{data.name}{\texttt{tactics-v1}}
\setvalue{data.rows}{810{,}140}
\setvalue{data.rows.test}{81{,}135}
\setvalue{data.share.test}{10}
\setvalue{data.share.train}{85}
\setvalue{data.share.valid}{5}
\setvalue{describe.mover-elo.q5}{1611}
\setvalue{describe.mover-elo.q50}{1876}
\setvalue{describe.solver-moves.q50}{2}
\setvalue{describe.solver-moves.q95}{4}
\setvalue{describe.win-before.min}{50}
\setvalue{describe.win-before.q50}{91.3}
\setvalue{describe.win-prior.q50}{42.1}
\setvalue{examples.count}{three}
\setvalue{floor.brier}{0.2145}
\setvalue{floor.constant}{69.0\%}
\setvalue{floor.pr-auc}{0.6884}
\setvalue{floor.score}{0.6203}
\setvalue{grid.ratings.-1700.-1000.happened}{83.9\%}
\setvalue{grid.ratings.-1700.1800-2200.happened}{49.8\%}
\setvalue{grid.ratings.-1700.2200+.happened}{46.2\%}
\setvalue{grid.ratings.-1700.2200+.said}{52.7\%}
\setvalue{grid.ratings.2300+.-1000.happened}{79.9\%}
\setvalue{grid.ratings.2300+.1800-2200.happened}{65.4\%}
\setvalue{grid.ratings.2300+.2200+.happened}{51.5\%}
\setvalue{grid.ratings.2300+.2200+.said}{63.2\%}
\setvalue{grid.ratings.gbm-hand.widest-off}{12.1}
\setvalue{grid.ratings.puzzle-rating.widest-off}{10.3}
\setvalue{grid.ratings.transformer-medium-hand-pretrained.widest-off}{11.7}
\setvalue{grid.ratings.transformer-medium-hand-pretrained.widest-off.column}{2200+}
\setvalue{grid.ratings.transformer-medium-hand-pretrained.widest-off.row}{2300+}
\setvalue{grid.smallest-cell}{30}
\setvalue{label.rate}{69.0}
\setvalue{label.rate.test}{68.8}
\setvalue{label.rate.train}{69.0}
\setvalue{label.rate.valid}{68.9}
\setvalue{ladder.at-last-epoch.of}{two of the four networks}
\setvalue{ladder.behind.of}{three of the four networks}
\setvalue{look.date}{23 November 2026}
\setvalue{look.resamples}{200}
\setvalue{look.smallest-compared}{200}
\setvalue{look.test-rows}{81{,}135}
\setvalue{manifest.dropped.no-clock}{66{,}930}
\setvalue{manifest.dropped.no-reply}{5{,}309}
\setvalue{manifest.passed-recipe}{882{,}379}
\setvalue{manifest.source-rows-read}{1{,}484{,}945}
\setvalue{method.gbm.features.percent}{80}
\setvalue{method.gbm.leaves}{127}
\setvalue{method.gbm.patience}{100}
\setvalue{method.gbm.rate}{0.05}
\setvalue{method.gbm.rounds}{4{,}000}
\setvalue{method.gbm.rows.percent}{80}
\setvalue{method.resamples}{1{,}000}
\setvalue{more.gbm-hand.most}{twice}
\setvalue{networks.ahead.of}{five of the nine networks}
\setvalue{networks.patience}{three}
\setvalue{networks.warmup.percent}{5}
\setvalue{pair.conv-medium-hand.transformer-medium-hand.gain}{+0.003182}
\setvalue{pair.conv-medium-hand.transformer-medium-hand.gain.high}{+0.004398}
\setvalue{pair.conv-medium-hand.transformer-medium-hand.gain.low}{+0.001855}
\setvalue{pair.conv-medium-hand.transformer-medium-hand.kept}{accepted}
\setvalue{pair.conv-medium-hand.transformer-medium-hand.valid.verdict}{better}
\setvalue{pair.conv-medium-hand.transformer-medium-hand.verdict}{significantly better}
\setvalue{pair.gbm-hand-double.gbm-hand.gain}{+0.003246}
\setvalue{pair.gbm-hand-double.gbm-hand.kept}{accepted}
\setvalue{pair.gbm-hand-double.gbm-hand.verdict}{significantly better}
\setvalue{pair.gbm-hand-half.gbm-hand.gain}{$-$0.003016}
\setvalue{pair.gbm-hand-half.gbm-hand.verdict}{significantly worse}
\setvalue{pair.gbm-hand-no-first-move.gbm-hand.gain}{$-$0.01330}
\setvalue{pair.gbm-hand-no-hanging.gbm-hand.gain}{$-$0.002901}
\setvalue{pair.gbm-hand-no-kings.gbm-hand.gain}{$-$0.0002682}
\setvalue{pair.gbm-hand-no-kings.gbm-hand.verdict}{not significantly different}
\setvalue{pair.gbm-hand-no-line.gbm-hand.gain}{$-$0.005165}
\setvalue{pair.gbm-hand-no-material.gbm-hand.gain}{$-$0.002310}
\setvalue{pair.gbm-hand-no-mobility.gbm-hand.gain}{$-$0.003835}
\setvalue{pair.gbm-hand-no-pawns.gbm-hand.gain}{$-$0.0001288}
\setvalue{pair.gbm-hand-no-pawns.gbm-hand.verdict}{not significantly different}
\setvalue{pair.gbm-hand-no-setup.gbm-hand.gain}{$-$0.002627}
\setvalue{pair.gbm-hand-quarter.gbm-hand.gain}{$-$0.007085}
\setvalue{pair.gbm-hand-quarter.gbm-hand.verdict}{significantly worse}
\setvalue{pair.hybrid-medium-hand.transformer-medium-hand.gain}{$-$0.002856}
\setvalue{pair.hybrid-medium-hand.transformer-medium-hand.gain.high}{$-$0.001651}
\setvalue{pair.hybrid-medium-hand.transformer-medium-hand.gain.low}{$-$0.004096}
\setvalue{pair.hybrid-medium-hand.transformer-medium-hand.kept}{not accepted}
\setvalue{pair.hybrid-medium-hand.transformer-medium-hand.verdict}{significantly worse}
\setvalue{pair.transformer-medium-attacks.transformer-medium.gain}{$-$0.0001147}
\setvalue{pair.transformer-medium-attacks.transformer-medium.gain.high}{+0.0004012}
\setvalue{pair.transformer-medium-attacks.transformer-medium.gain.low}{$-$0.0006363}
\setvalue{pair.transformer-medium-attacks.transformer-medium.kept}{not accepted}
\setvalue{pair.transformer-medium-attacks.transformer-medium.verdict}{not significantly different}
\setvalue{pair.transformer-medium-hand-attacks.transformer-medium-hand.gain}{$-$0.0000462}
\setvalue{pair.transformer-medium-hand-attacks.transformer-medium-hand.gain.high}{+0.0000439}
\setvalue{pair.transformer-medium-hand-attacks.transformer-medium-hand.gain.low}{$-$0.0001389}
\setvalue{pair.transformer-medium-hand-attacks.transformer-medium-hand.kept}{not accepted}
\setvalue{pair.transformer-medium-hand-attacks.transformer-medium-hand.verdict}{not significantly different}
\setvalue{pair.transformer-medium-hand-double.transformer-medium-hand.gain}{+0.002131}
\setvalue{pair.transformer-medium-hand-double.transformer-medium-hand.kept}{accepted}
\setvalue{pair.transformer-medium-hand-double.transformer-medium-hand.verdict}{significantly better}
\setvalue{pair.transformer-medium-hand-half.transformer-medium-hand.gain}{$-$0.007087}
\setvalue{pair.transformer-medium-hand-long.transformer-medium-hand.gain}{$-$0.000357}
\setvalue{pair.transformer-medium-hand-long.transformer-medium-hand.gain.high}{+0.000428}
\setvalue{pair.transformer-medium-hand-long.transformer-medium-hand.gain.low}{$-$0.001175}
\setvalue{pair.transformer-medium-hand-long.transformer-medium-hand.kept}{not accepted}
\setvalue{pair.transformer-medium-hand-long.transformer-medium-hand.valid.verdict}{better}
\setvalue{pair.transformer-medium-hand-long.transformer-medium-hand.verdict}{not significantly different}
\setvalue{pair.transformer-medium-hand-no-geometry.transformer-medium-hand.gain}{$-$0.0001291}
\setvalue{pair.transformer-medium-hand-no-geometry.transformer-medium-hand.gain.high}{+0.0000451}
\setvalue{pair.transformer-medium-hand-no-geometry.transformer-medium-hand.gain.low}{$-$0.0002895}
\setvalue{pair.transformer-medium-hand-no-geometry.transformer-medium-hand.verdict}{not significantly different}
\setvalue{pair.transformer-medium-hand-pretrained.transformer-medium-hand.gain}{+0.008191}
\setvalue{pair.transformer-medium-hand-pretrained.transformer-medium-hand.gain.high}{+0.009645}
\setvalue{pair.transformer-medium-hand-pretrained.transformer-medium-hand.gain.low}{+0.006677}
\setvalue{pair.transformer-medium-hand-pretrained.transformer-medium-hand.kept}{accepted}
\setvalue{pair.transformer-medium-hand-pretrained.transformer-medium-hand.valid.verdict}{better}
\setvalue{pair.transformer-medium-hand-pretrained.transformer-medium-hand.verdict}{significantly better}
\setvalue{pair.transformer-medium-hand-quarter.transformer-medium-hand.gain}{$-$0.01490}
\setvalue{pair.transformer-medium-hand.transformer-medium.gain}{+0.01481}
\setvalue{pair.transformer-medium-hand.transformer-medium.gain.high}{+0.01628}
\setvalue{pair.transformer-medium-hand.transformer-medium.gain.low}{+0.01332}
\setvalue{pair.transformer-medium-hand.transformer-medium.kept}{accepted}
\setvalue{pair.transformer-medium-hand.transformer-medium.valid.verdict}{better}
\setvalue{pair.transformer-medium-hand.transformer-medium.verdict}{significantly better}
\setvalue{pair.transformer-medium-long.transformer-medium.gain}{+0.001499}
\setvalue{pair.transformer-medium-long.transformer-medium.gain.high}{+0.002272}
\setvalue{pair.transformer-medium-long.transformer-medium.gain.low}{+0.000711}
\setvalue{pair.transformer-medium-long.transformer-medium.kept}{accepted}
\setvalue{pair.transformer-medium-long.transformer-medium.valid.verdict}{better}
\setvalue{pair.transformer-medium-long.transformer-medium.verdict}{significantly better}
\setvalue{pair.transformer-medium-pretrained.transformer-medium.gain}{+0.01887}
\setvalue{pair.transformer-medium-pretrained.transformer-medium.gain.high}{+0.02053}
\setvalue{pair.transformer-medium-pretrained.transformer-medium.gain.low}{+0.01719}
\setvalue{pair.transformer-medium-pretrained.transformer-medium.kept}{accepted}
\setvalue{pair.transformer-medium-pretrained.transformer-medium.valid.verdict}{better}
\setvalue{pair.transformer-medium-pretrained.transformer-medium.verdict}{significantly better}
\setvalue{recipe.files}{7}
\setvalue{recipe.filters.max-rating-deviation}{100}
\setvalue{recipe.filters.min-plays}{100}
\setvalue{reference.bottom.10.rate}{66.1}
\setvalue{reference.brier}{0.1828}
\setvalue{reference.calibration-error}{0.0043}
\setvalue{reference.calibration.highest}{94.8}
\setvalue{reference.calibration.lowest}{33.8}
\setvalue{reference.calibration.widest-gap}{0.9}
\setvalue{reference.features}{47}
\setvalue{reference.leans.1.feature}{win probability before the move}
\setvalue{reference.leans.1.share}{9}
\setvalue{reference.leans.2.feature}{first move is a sacrifice}
\setvalue{reference.leans.2.share}{9}
\setvalue{reference.leans.3.feature}{win probability before the opponent's mistake}
\setvalue{reference.leans.3.share}{6}
\setvalue{reference.leans.4.feature}{mover's rating}
\setvalue{reference.leans.4.share}{5}
\setvalue{reference.leans.5.feature}{mover's clock}
\setvalue{reference.leans.5.share}{5}
\setvalue{reference.platt.intercept}{$-$0.02}
\setvalue{reference.platt.slope}{1.02}
\setvalue{reference.pr-auc}{0.8560}
\setvalue{reference.roc-auc}{0.7337}
\setvalue{reference.score}{0.5414}
\setvalue{reference.trees}{787}
\setvalue{results.date}{23 November 2026}
\setvalue{run.conv-medium-hand.lr}{0.002}
\setvalue{run.conv-medium-hand.parameters.millions}{7.0}
\setvalue{run.conv-medium-hand.score}{0.5337}
\setvalue{run.conv-small.best-epoch}{5}
\setvalue{run.conv-small.epochs}{8}
\setvalue{run.conv-small.gain}{$-$0.00997}
\setvalue{run.conv-small.gain.high}{$-$0.00810}
\setvalue{run.conv-small.gain.low}{$-$0.01184}
\setvalue{run.conv-small.parameters.millions}{1.0}
\setvalue{run.conv-small.score}{0.5513}
\setvalue{run.gbm-context.pr-auc}{0.7961}
\setvalue{run.gbm-context.roc-auc}{0.6369}
\setvalue{run.gbm-context.score}{0.5923}
\setvalue{run.gbm-hand-double.score}{0.5381}
\setvalue{run.gbm-hand-double.trees}{1{,}202}
\setvalue{run.gbm-hand-half.score}{0.5444}
\setvalue{run.gbm-hand-no-first-move.score}{0.5547}
\setvalue{run.gbm-hand-quarter.score}{0.5485}
\setvalue{run.hybrid-medium-hand.convs}{2}
\setvalue{run.hybrid-medium-hand.depth}{5}
\setvalue{run.hybrid-medium-hand.gain}{+0.001587}
\setvalue{run.hybrid-medium-hand.gain.high}{+0.002928}
\setvalue{run.hybrid-medium-hand.gain.low}{+0.000300}
\setvalue{run.hybrid-medium-hand.parameters.millions}{6.6}
\setvalue{run.puzzle-rating.calibration-error}{0.0171}
\setvalue{run.puzzle-rating.gain}{$-$0.04753}
\setvalue{run.puzzle-rating.gain.high}{$-$0.04484}
\setvalue{run.puzzle-rating.gain.low}{$-$0.04996}
\setvalue{run.puzzle-rating.platt.slope}{0.19}
\setvalue{run.puzzle-rating.pr-auc}{0.7869}
\setvalue{run.puzzle-rating.roc-auc}{0.6570}
\setvalue{run.puzzle-rating.score}{0.5889}
\setvalue{run.transformer-medium-attacks.gain}{$-$0.01048}
\setvalue{run.transformer-medium-attacks.gain.high}{$-$0.00881}
\setvalue{run.transformer-medium-attacks.gain.low}{$-$0.01211}
\setvalue{run.transformer-medium-attacks.score}{0.5518}
\setvalue{run.transformer-medium-hand-double.best-epoch}{6}
\setvalue{run.transformer-medium-hand-double.epochs}{6}
\setvalue{run.transformer-medium-hand-double.score}{0.5348}
\setvalue{run.transformer-medium-hand-half.score}{0.5440}
\setvalue{run.transformer-medium-hand-long.best-epoch}{13}
\setvalue{run.transformer-medium-hand-pretrained.best-epoch}{8}
\setvalue{run.transformer-medium-hand-pretrained.epochs}{12}
\setvalue{run.transformer-medium-hand-pretrained.over-twin}{+0.00413}
\setvalue{run.transformer-medium-hand-pretrained.score}{0.5287}
\setvalue{run.transformer-medium-hand-quarter.score}{0.5518}
\setvalue{run.transformer-medium-hand.best-epoch}{10}
\setvalue{run.transformer-medium-hand.epochs}{12}
\setvalue{run.transformer-medium-hand.gain}{+0.004443}
\setvalue{run.transformer-medium-hand.gain.high}{+0.005730}
\setvalue{run.transformer-medium-hand.gain.low}{+0.003247}
\setvalue{run.transformer-medium-hand.lr}{0.0006}
\setvalue{run.transformer-medium-hand.parameters.millions}{6.6}
\setvalue{run.transformer-medium-hand.platt.slope}{0.90}
\setvalue{run.transformer-medium-hand.score}{0.5369}
\setvalue{run.transformer-medium-long.best-epoch}{15}
\setvalue{run.transformer-medium-long.epochs}{24}
\setvalue{run.transformer-medium-long.gain}{$-$0.00887}
\setvalue{run.transformer-medium-long.gain.high}{$-$0.00719}
\setvalue{run.transformer-medium-long.gain.low}{$-$0.01048}
\setvalue{run.transformer-medium-pretrained.best-epoch}{8}
\setvalue{run.transformer-medium-pretrained.epochs}{12}
\setvalue{run.transformer-medium-pretrained.gain}{+0.00851}
\setvalue{run.transformer-medium-pretrained.gain.high}{+0.01035}
\setvalue{run.transformer-medium-pretrained.gain.low}{+0.00671}
\setvalue{run.transformer-medium-pretrained.score}{0.5329}
\setvalue{run.transformer-medium.best-epoch}{12}
\setvalue{run.transformer-medium.epochs}{12}
\setvalue{run.transformer-medium.gain}{$-$0.01037}
\setvalue{run.transformer-medium.gain.high}{$-$0.00868}
\setvalue{run.transformer-medium.gain.low}{$-$0.01199}
\setvalue{run.transformer-medium.parameters.millions}{6.6}
\setvalue{run.transformer-medium.score}{0.5517}
\setvalue{run.transformer-small.best-epoch}{8}
\setvalue{run.transformer-small.epochs}{8}
\setvalue{run.transformer-small.gain}{$-$0.01175}
\setvalue{run.transformer-small.gain.high}{$-$0.01015}
\setvalue{run.transformer-small.gain.low}{$-$0.01340}
\setvalue{run.transformer-small.parameters.millions}{0.9}
\setvalue{run.transformer-small.score}{0.5531}
\setvalue{runs.count}{14}
\setvalue{slice.backward.no.happened}{68.8\%}
\setvalue{slice.backward.yes.happened}{69.1\%}
\setvalue{slice.band.-1700.happened}{66.0\%}
\setvalue{slice.band.2100-2299.happened}{71.5\%}
\setvalue{slice.band.ahead.of}{all five groups}
\setvalue{slice.band.challenger.widest-off}{1.1}
\setvalue{slice.band.reference.widest-off}{0.8}
\setvalue{slice.capture.no.happened}{66.5\%}
\setvalue{slice.capture.yes.happened}{72.0\%}
\setvalue{slice.check.no.gain}{+0.01811}
\setvalue{slice.check.no.gain.high}{+0.02066}
\setvalue{slice.check.no.gain.low}{+0.01590}
\setvalue{slice.check.no.happened}{65.2\%}
\setvalue{slice.check.yes.gain}{+0.007945}
\setvalue{slice.check.yes.gain.high}{+0.009972}
\setvalue{slice.check.yes.gain.low}{+0.005341}
\setvalue{slice.check.yes.happened}{72.0\%}
\setvalue{slice.clock.-10.challenger.off}{+4.6}
\setvalue{slice.clock.-10.gain}{$-$0.00040}
\setvalue{slice.clock.-10.happened}{58.0\%}
\setvalue{slice.clock.-10.reference.off}{+1.0}
\setvalue{slice.clock.-10.share}{2}
\setvalue{slice.clock.10-30.challenger.off}{+2.7}
\setvalue{slice.clock.10-30.happened}{65.1\%}
\setvalue{slice.clock.10-30.reference.off}{$-$0.3}
\setvalue{slice.clock.10-30.share}{7}
\setvalue{slice.clock.180-600.happened}{69.1\%}
\setvalue{slice.clock.30-60.happened}{68.9\%}
\setvalue{slice.clock.60-180.happened}{68.7\%}
\setvalue{slice.clock.600+.happened}{72.4\%}
\setvalue{slice.clock.ahead.groups}{10-30, 30-60, 60-180, 180-600 and 600+}
\setvalue{slice.clock.ahead.of}{five of the six groups}
\setvalue{slice.clock.challenger.widest-off}{4.6}
\setvalue{slice.first-piece.challenger.widest-off}{3.0}
\setvalue{slice.first-piece.challenger.widest-off.group}{king}
\setvalue{slice.first-piece.king.happened}{75.1\%}
\setvalue{slice.first-piece.pawn.happened}{66.5\%}
\setvalue{slice.mate.no.happened}{66.0\%}
\setvalue{slice.mate.yes.happened}{79.3\%}
\setvalue{slice.moves.1.gain}{+0.002801}
\setvalue{slice.moves.1.happened}{86.7\%}
\setvalue{slice.moves.2.happened}{71.2\%}
\setvalue{slice.moves.3.happened}{61.2\%}
\setvalue{slice.moves.5+.happened}{59.0\%}
\setvalue{slice.prior.10-30.happened}{66.9\%}
\setvalue{slice.prior.30-50.happened}{65.3\%}
\setvalue{slice.prior.90+.happened}{90.2\%}
\setvalue{slice.prior.90+.share}{5}
\setvalue{slice.prior.ahead.groups}{\textless{}10, 10-30, 30-50 and 50-70}
\setvalue{slice.prior.ahead.of}{four of the six groups}
\setvalue{slice.prior.challenger.widest-off}{1.6}
\setvalue{slice.prior.reference.widest-off}{0.6}
\setvalue{slice.puzzle-rating.-1000.challenger.off}{$-$2.3}
\setvalue{slice.puzzle-rating.-1000.happened}{84.0\%}
\setvalue{slice.puzzle-rating.-1000.reference.off}{$-$3.5}
\setvalue{slice.puzzle-rating.2200+.challenger.off}{+9.7}
\setvalue{slice.puzzle-rating.2200+.happened}{48.7\%}
\setvalue{slice.puzzle-rating.2200+.reference.off}{+10.3}
\setvalue{slice.puzzle-rating.ahead.of}{all eight groups}
\setvalue{slice.puzzle-rating.challenger.widest-off}{9.7}
\setvalue{slice.puzzle-rating.gain.most}{+0.01911}
\setvalue{slice.puzzle-rating.gain.most.group}{2200+}
\setvalue{slice.quiet.0.gain}{+0.01062}
\setvalue{slice.quiet.0.happened}{72.0\%}
\setvalue{slice.quiet.1.happened}{63.3\%}
\setvalue{slice.quiet.2.gain}{+0.02263}
\setvalue{slice.sacrifice.no.happened}{72.0\%}
\setvalue{slice.sacrifice.yes.happened}{51.0\%}
\setvalue{slice.sacrifice.yes.share}{15}
\setvalue{slice.seconds.-2.happened}{75.3\%}
\setvalue{slice.seconds.2-5.challenger.off}{$-$4.0}
\setvalue{slice.seconds.2-5.happened}{75.9\%}
\setvalue{slice.seconds.2-5.reference.off}{$-$4.7}
\setvalue{slice.seconds.40+.challenger.off}{+10.3}
\setvalue{slice.seconds.40+.happened}{51.5\%}
\setvalue{slice.seconds.40+.reference.off}{+12.7}
\setvalue{slice.speed.blitz.challenger.off}{+1.3}
\setvalue{slice.speed.blitz.happened}{66.2\%}
\setvalue{slice.speed.classical.challenger.off}{$-$1.2}
\setvalue{slice.speed.classical.happened}{73.7\%}
\setvalue{slice.speed.rapid.happened}{69.4\%}
\setvalue{slice.themes.ahead.of}{22 of the 30 groups}
\setvalue{slice.themes.attraction.happened}{51.7\%}
\setvalue{slice.themes.backRankMate.happened}{82.6\%}
\setvalue{slice.themes.hangingPiece.happened}{80.7\%}
\setvalue{slice.themes.mateIn1.happened}{88.5\%}
\setvalue{slice.themes.oneMove.happened}{86.7\%}
\setvalue{slice.themes.quietMove.gain}{+0.02710}
\setvalue{slice.themes.reference.widest-off}{3.2}
\setvalue{slice.themes.reference.widest-off.group}{sacrifice}
\setvalue{slice.themes.sacrifice.happened}{48.2\%}
\setvalue{study.features.gbm-hand.worse.of}{six of the eight ablations}
\setvalue{study.games.transformer-medium-hand.worse.of}{both ablations}
\setvalue{task.context.count}{14}
\setvalue{task.hand.count}{33}
\setvalue{task.marks.count}{14}
\setvalue{task.reference}{\run{gbm-hand}}


\begin{document}

\papertitle{Predicting Whether Players Find Tactics in Their Own Games}
  {Results as of \V{results.date}\quad\textperiodcentered\quad dataset \V{data.name},
   build \V{data.fingerprint}\quad\textperiodcentered\quad code \V{code.commit}}

\begin{abstract}
A chess puzzle tells its solver that a tactic exists. In a game the player
has to notice it, with the clock running. We train models that predict
whether a player finds a tactic that the opponent's last move has made
possible. The inputs are the position, the opponent's move, the solution of
the tactic, both players' ratings, both clocks and the engine evaluations
before and after the opponent's move. The dataset joins \V{data.rows} Lichess
puzzles to the rated games from which they were generated, so that the next
move of the game shows whether the tactic was found, and it is split by game.
Players found the tactic in \V{label.rate}\% of the positions. We compare a
constant baseline, gradient-boosted decision trees (GBDT) on hand-crafted
features, a baseline that applies the Elo formula to the player's rating and
the puzzle's rating, and convolutional and transformer networks that read the
board as 64 square tokens. The best of \V{runs.count} models is a transformer
that also receives the hand-crafted features and whose encoder is pretrained
on engine evaluations. On \V{data.rows.test} test positions its log loss is
\V{best.score} (95\% CI \V{best.score.low} to \V{best.score.high}), compared
with \V{reference.score} for the GBDT baseline, \V{run.puzzle-rating.score}
for the rating-based baseline and \V{floor.score} for the constant baseline.
Trained from random initialisation without the hand-crafted features, a small
residual network and transformers of two sizes are all significantly worse
than the GBDT baseline. The hand-crafted features, pretraining, a residual
convolutional encoder and more training data each improve a network
significantly. After calibration, the predicted probabilities are within
\V{best.calibration.widest-gap} percentage points of the observed rate in
each of \V{best.calibration.groups} bins\iflook. For the tactics with the
highest puzzle ratings, however, the mean prediction of the best model is
\V{slice.puzzle-rating.challenger.widest-off} percentage points above the
observed share\fi.
\end{abstract}

\section{Introduction}
\label{sec:intro}

When a game is reviewed with a chess engine, one situation recurs. The
opponent made a mistake, a tactic was available, and the player did not play
it. How much attention this deserves depends on how hard the tactic was to
find in that game. If players of the same strength almost always find it,
the miss was a lapse. If most of them miss it too, the tactic was beyond
what can be expected at that level, and pointing it out teaches little.

The rating of a puzzle does not answer this question. Lichess rates each of
its puzzles with the Glicko-2 system~\cite{glicko2} from the attempts of its
solvers, so the rating measures how often a tactic is solved by players who
know that there is one. In a game nobody says so. The player has to notice
that the opponent's last move changed the position, while the clock is
running.

The data needed to study the question directly are public. Lichess generates
its puzzles from games played on the site~\cite{lichessdb}, and a public
dataset joins each puzzle to the game it comes from~\cite{puzzlegames}. A
puzzle begins after a move by the opponent, and its solution begins with the
reply that exploits this move. The record of the game contains the move that
the player actually made at this point, together with both ratings, the clock
times and the engine evaluations. Every puzzle is therefore also a record of
whether one player, with a known rating and a known amount of time, found the
tactic in a game.

We study the following prediction task. Given the position after the
opponent's move, the solution of the tactic, both players' ratings, both
clock times and the engine evaluations before and after the opponent's move,
predict the probability that the move played is the first move of the
solution. The model is given the tactic. It does not have to find it, which
is the task of an engine. It has to estimate how likely a human player in
this situation is to play it.

This work supports a chess coaching application that would use the
prediction in a game review. A missed tactic with a high predicted
probability is a lapse, because most players of the user's strength find it.
A missed tactic with a low predicted probability is one that most of them
miss as well. The two cases call for different feedback.

We evaluate a fixed sequence of models on \V{data.rows} positions from rated
games on Lichess: a constant baseline, GBDT on scalar context features, GBDT
with hand-crafted features of the position and of the solution, and neural
networks that read the board directly, with and without the hand-crafted
features. Because the sequence is fixed, each step measures the contribution
of one component. Our main findings are the following.
\begin{itemize}
\item \textbf{The description of the tactic matters more than the context.}
GBDT on the ratings, the clocks, the two win probabilities and the length of
the solution, without the board, has a log loss of
\V{run.gbm-context.score}. It achieves \V{context.lift.share}\% of the
improvement of the GBDT baseline (\V{reference.score}) over the constant
baseline (\V{floor.score}). The larger part comes from the hand-crafted
features, which describe the position and the moves of the solution.
\item \textbf{The puzzle's rating and the mover's rating predict poorly on
their own.} The Elo formula applied to the two ratings, with a slope and an
intercept fitted on the validation set, has a log loss of
\V{run.puzzle-rating.score}. This is close to that of GBDT on the context
features and significantly worse than the GBDT baseline. The fitted slope is
\V{run.puzzle-rating.platt.slope}, where a slope of one would mean that a
rating difference has the effect in a game that the formula assigns to it
(Section~\ref{sec:rating-baseline}).
\item \textbf{From random initialisation, a network needs the hand-crafted
features.} Without them, \V{ladder.behind.of} of the fixed sequence are
significantly worse than the GBDT baseline. With them, the medium transformer
is significantly better (gain \V{run.transformer-medium-hand.gain}; 95\% CI
\V{run.transformer-medium-hand.gain.low} to
\V{run.transformer-medium-hand.gain.high}).
\ifchangestart
\item \textbf{Pretraining is the most effective modification.} Initialising
the transformer's encoder from a model trained on engine evaluations gives a
paired gain of \V{pair.transformer-medium-pretrained.transformer-medium.gain}
in log loss without the hand-crafted features and of
\V{pair.transformer-medium-hand-pretrained.transformer-medium-hand.gain} with
them.\ifchangeconv{} A residual convolutional encoder in place of the
transformer
(\V{pair.conv-medium-hand.transformer-medium-hand.gain})\ifmore{} and twice
the training data
(\V{pair.transformer-medium-hand-double.transformer-medium-hand.gain})\fi{}
are also accepted\ifchangeattacks, and an attack-aware attention bias is
not\fi{} (Section~\ref{sec:ablations}).\fi
\fi
\item \textbf{The predicted probabilities are calibrated.} After calibration,
the mean predicted probability of the best model and the observed share of
tactics found differ by at most \V{best.calibration.widest-gap} percentage
points in each of \V{best.calibration.groups} bins\iflook. Among the 10\% of
test positions with the highest predictions, \V{best.top.10.rate}\% of the
tactics were found. Among the 10\% with the lowest predictions,
\V{best.bottom.10.rate}\% were missed, \V{best.bottom.10.lift} times the
share missed in the whole test set\fi.
\iflook
\item \textbf{Sacrifices, long solutions and long thinks go with missed
tactics.} The tactic is found in \V{slice.sacrifice.yes.happened} of the
positions in which its first move is a sacrifice, compared with
\V{slice.sacrifice.no.happened} of the others. The share found falls from
\V{slice.moves.1.happened} for solutions of one move to
\V{slice.moves.5+.happened} for solutions of five moves or more, and from
\V{slice.seconds.-2.happened} for moves played in less than two seconds to
\V{slice.seconds.40+.happened} for moves that took 40 seconds or more
(Section~\ref{sec:error}).
\item \textbf{The models overestimate how often the hardest tactics are
found.} For puzzles rated 2200 and above, the mean prediction of the best
model is \V{slice.puzzle-rating.challenger.widest-off} percentage points
above the observed share. The puzzle's
rating therefore contains information about a tactic that the models do not
extract from the position and the solution.
\fi
\end{itemize}

\section{Related Work}
\label{sec:related}

\paragraph{Human error in chess.}
Anderson et al.~\cite{humanerror} label errors exactly using endgame
tablebases and find that the difficulty of the position predicts errors
better than the player's skill or the time available. The Maia
study~\cite{maia} includes classifiers that predict, from the position and
from metadata about the players and the game, whether the next move will be a
blunder. Our task is narrower. Every position contains a tactic that Lichess
has published as a puzzle, and the question is whether the player finds it.

\paragraph{Puzzle difficulty.}
Lichess rates its puzzles with Glicko-2~\cite{glicko2}. Predicting this
rating from the position and the solution was the task of data-mining
competitions at IEEE BigData 2024 and FedCSIS
2025~\cite{puzzlecup,puzzlecupwinner}, and Mi{\l}osz and
Kapusta~\cite{glickformer} predict it with a transformer over the positions
of the solution. These studies predict how hard a tactic is for solvers of
puzzles. We predict whether it is found in a game. We use the puzzle's rating
only in one baseline and in the error analysis.

\paragraph{Models of human moves.}
Maia~\cite{maia} and Maia-2~\cite{maia2} predict the move that a player of a
given rating plays. The probability that such a model assigns to the first
move of the solution is a natural estimate of the quantity we study, and it
does not require the solution as an input. We do not evaluate it here
(Section~\ref{sec:limits}).

\section{Task and Data}
\label{sec:data}

\subsection{Task}

An example is one position of a rated game, reached by a move of the
opponent that makes a tactic possible. We call the player to move the
\emph{mover}, the opponent's preceding move the \emph{mistake}, and the
winning line of the corresponding Lichess puzzle the \emph{solution}. The
tables of input planes and features, which this task shares with a task on
puzzle ratings, call the mover the \emph{solver}. The input is the position,
the mistake, the solution, both players' ratings, both clock times, the time
control, the ply, and the mover's win probability before and after the
mistake. The target is binary: whether the mover found the tactic, that is
whether the move played in the game is the first move of the solution.
Table~\ref{tab:brief} summarises the task.

\begin{table}[htbp]\centering\small
\begin{tabularx}{\textwidth}{lL}
\toprule
Input & The position after the opponent's mistake, the mistake, the solution
of the tactic, both players' ratings, both clock times, the time control, the
ply, the mover's win probability before and after the mistake, and the
length of the solution. \\
Target & Whether the move played is the first move of the solution. \\
Data & Lichess puzzles joined to the rated games from which they were
generated (\V{data.licence}): blitz, rapid and classical games, and puzzles
whose rating is established. \\
Split & By game: \V{data.share.train}\% training, \V{data.share.valid}\%
validation, \V{data.share.test}\% test. \\
Primary metric & Log loss on the test set (lower is better). \\
\bottomrule
\end{tabularx}
\caption{Summary of the task.}
\label{tab:brief}
\end{table}

\paragraph{Definition of the label.}
A tactic counts as \emph{found} when the move played in the game is the first
move of the puzzle's solution. Lichess accepts any checkmate where the
solution of a puzzle is a mate in one, and we follow it: when the first move
of the solution mates, any move that mates counts as well. Only the first
move is compared. The label does not record whether the player went on to
play the rest of the solution.

Two consequences of this definition matter for the results. First, a mover
who plays a different move is counted as having missed the tactic, even if
that move also keeps the advantage. The dataset built for this paper does not
contain the evaluation after the move played, so we cannot say how often this
happens. Second, the label says nothing about why a move was played. A mover
who plays the first move of the solution for another reason, for example as
an obvious recapture, is counted as having found the tactic.

\paragraph{Win probability.}
Engine evaluations enter as win probabilities. As in the Lichess accuracy
metric~\cite{lichessaccuracy}, an evaluation of $c$ centipawns, limited to
the range from $-1000$ to $1000$ with a forced mate counted as the
corresponding limit, is mapped to a win probability in percent,
\[
W(c) = 50 + 50\left(\frac{2}{1 + e^{-0.00368208\,c}} - 1\right).
\]
The evaluations are those that Lichess attaches to the record of a game after
computer analysis. A model receives two of them, from the mover's point of
view: the win probability before the opponent's mistake and the win
probability after it, which is the evaluation of the position on the board.

\subsection{Dataset}

The source is a dataset published by Lichess that joins each puzzle of the
Lichess puzzle database~\cite{lichessdb} to the game from which the puzzle
was generated~\cite{puzzlegames}. For each puzzle it gives the position, the
solution, the puzzle's rating and theme tags, and the record of the game with
both ratings, the time control, the engine evaluation after every move and,
for most games, the clock time after every move. The source consists of fourteen files ordered by puzzle
identifier. The identifiers are random strings, so the first
\V{recipe.files} files, which we read, are in effect a random half of the
puzzles. Because the source may still change, a build reads one fixed
revision of it (Section~\ref{sec:cost}).

A puzzle is included when its rating is established, which we take to mean a
rating deviation of at most \V{recipe.filters.max-rating-deviation} after at
least \V{recipe.filters.min-plays} attempts, and when its game is a game of
standard chess at blitz, rapid or classical speed in which both players have
a rating and neither is a bot. Bullet games are excluded, because at that
speed many moves are entered before the opponent's move has been seen. Of
\V{manifest.source-rows-read} puzzles read, \V{manifest.passed-recipe} meet
these conditions. Of these, \V{manifest.dropped.no-clock} are dropped because
the record of the game has no clock times and \V{manifest.dropped.no-reply}
because the game ended with the opponent's mistake, which leaves
\V{data.rows} positions (Table~\ref{tab:data}).

For each puzzle, the example is the position after the opponent's mistake at
the ply that the puzzle names, and the label compares the next move of the
game with the solution. The mover's clock is the time remaining after the
mover's previous move, and the opponent's clock is the time remaining after
the mistake. The mover found the tactic in \V{label.rate}\% of the positions
(\V{label.rate.train}\% in the training set, \V{label.rate.valid}\% in the
validation set and \V{label.rate.test}\% in the test set).

The split is determined by a hash of the game identifier, and the same
grouping is used for the bootstrap (Section~\ref{sec:protocol}). In the
training and test sets every game contributes exactly one position, so the
\V{data.rows.test} test positions come from as many games. The dataset is not
stored in the repository. Its recipe is, and a build of the dataset is
identified by a fingerprint, a hash of its contents. All models in this paper
were trained and evaluated on build \V{data.fingerprint}\ifmore; the models
of Section~\ref{sec:training-size} that use more training games were trained
on a larger build, which reads all fourteen files, and evaluated on the
validation and test sets of this one\fi. The evaluation code refuses to rank
models evaluated on different builds.

\begin{table}[htbp]\centering\small
\block{data}
\caption{Dataset statistics by split.}
\label{tab:data}
\end{table}

\iflook
Table~\ref{tab:describe} summarises the main variables. Three properties of
the dataset should be kept in mind. First, the movers are strong: their
median rating is \V{describe.mover-elo.q50}, and only 5\% of them are rated
below \V{describe.mover-elo.q5}. Second, the positions are favourable to the
mover by construction. The median win probability is
\V{describe.win-prior.q50}\% before the mistake and
\V{describe.win-before.q50}\% after it, and it is at least
\V{describe.win-before.min}\% after it in every position. Third, most
tactics are short: the median solution has \V{describe.solver-moves.q50}
moves by the solver, and 95\% have at most \V{describe.solver-moves.q95}.

\begin{table}[htbp]\centering\small
\block{describe}
\caption{Summary statistics of selected variables over the whole dataset:
the mean, five quantiles, the minimum and the maximum. The puzzle's rating
and the time spent on the move are not inputs of the models of the fixed
sequence.}
\label{tab:describe}
\end{table}
\fi

\section{Methods}
\label{sec:methods}

\subsection{Input representation}
\label{sec:input}

% Input representation: identical for every task, so each model paper inputs
% this file under its own heading and then adds what is specific to its task.
% The prose follows docs/writing.md.
All tasks in this project share one input representation, so two tasks differ
in their target and in the information a model is given, and not in how a
position is encoded.

\paragraph{Square tokens.}
A position is encoded as 64 tokens, one per square. Token 0 denotes an empty
square, tokens 1 to 6 the pawn, knight, bishop, rook, queen and king of the
side to move, which we call the \emph{mover}, and tokens 7 to 12 the
opponent's pieces. When Black is to move, the board is flipped vertically and
the colours are swapped, so that the mover always plays up the board. A
position with White to move and the corresponding colour-reversed position
with Black to move therefore have the same encoding, and a model does not
have to learn the two cases separately.

\paragraph{Auxiliary input planes.}
Information about individual squares is supplied as binary planes over the 64
squares, which a network receives together with the tokens. This task uses
\V{task.marks.count} planes (Table~\ref{tab:marks}). The first marks the en
passant target square, and the others mark the origin and destination squares
of moves. A plane identifies the squares of a move and carries no other
information about it, such as whether the move is a capture or a check.

\begin{table}[htbp]\centering\small
\block{marks}
\caption{Auxiliary input planes. Squares are given from the mover's point of
view, after the board has been flipped.}
\label{tab:marks}
\end{table}

\ifsequence
\paragraph{Sequences of positions.}
An example of this task is not a single position. It is the sequence of one
player's moves in a game, and each element of the sequence is a position
encoded as above: the board before the move, with the opponent's preceding
move and the move played marked in the auxiliary planes. Scalar inputs
therefore come at two levels, one for each move and one for the game.

\paragraph{Per-move features.}
Each position of the sequence has \V{task.steps.count} scalar features of its
own (Table~\ref{tab:steps}): the four castling rights and the task's
information about that move. A network receives them together with the
position they belong to.

\begin{table}[htbp]\centering\small
\block{steps}
\caption{Per-move features: the scalar inputs attached to each position of a
sequence.}
\label{tab:steps}
\end{table}

\paragraph{Scalar context features.}
Information about the game as a whole is given as \V{task.context.count}
scalar context features (Table~\ref{tab:reads}). Every model receives them,
including the gradient-boosted trees, which never see a board. For a network,
both kinds of scalar feature are standardised with the mean and standard
deviation of the training set.

\begin{table}[htbp]\centering\small
\block{reads}
\caption{Scalar context features: the game-level inputs given to every
model.}
\label{tab:reads}
\end{table}

\paragraph{Hand-crafted features.}
Gradient-boosted trees cannot take a sequence of boards as input, so we also
compute \V{task.hand.count} aggregate features over the player's moves
(Table~\ref{tab:hand} in the appendix): how much win probability the moves
lost on average, at worst and how often beyond fixed thresholds, how the clock
was used, and how many moves were available. These features have two roles.
They are the inputs of the GBDT baseline, and so define the level that a
network reading every position has to exceed. They can also be given to a
network as additional context features, which tests whether the network
already extracts the same information from the sequence.

\else
\paragraph{Scalar context features.}
Information that does not belong to a square is given as scalar context
features: the four castling rights and the scalar inputs specific to the
task. This task uses \V{task.context.count} (Table~\ref{tab:reads}). Every
model receives them, including the gradient-boosted trees, which never see
the board. For a network they are standardised with the mean and standard
deviation of the training set.

\begin{table}[htbp]\centering\small
\block{reads}
\caption{Scalar context features: the inputs given to every model in addition
to the board.}
\label{tab:reads}
\end{table}

\paragraph{Hand-crafted features.}
Gradient-boosted trees cannot take 64 tokens as input, so we also compute
\V{task.hand.count} hand-crafted features of the position
(Table~\ref{tab:hand} in the appendix): material, mobility, available checks
and captures, hanging pieces, passed pawns and attacks near each king, and,
where the task gives a line of moves, features of that line. These features
have two roles. They are the inputs of the GBDT baseline, and so define the
level that a network reading the board has to exceed. They can also be given
to a network as additional context features, which tests whether the network
already extracts the same information from the board.

\fi

\iflook
\begin{figure}[htbp]\centering
\block{encoding}
\ifsequence
\caption{One position of one example, shown twice: the position before the
player's most costly move of the game. Left: as played, with arrows for the
opponent's preceding move (grey) and the player's move (blue). Right: as
encoded, oriented so that the mover plays up the board, with the squares set
in an auxiliary plane shaded grey for the opponent's move and blue for the
player's move. Every position of the sequence is encoded in this way.}
\else
\caption{One example from the dataset, shown twice. Left: as played, with
arrows for the opponent's preceding move (grey) and the move the example
concerns (blue). Right: as encoded, oriented so that the mover plays up the
board, with the squares set in an auxiliary plane shaded grey for the
opponent's move and, where the task marks a line of moves, blue for the
mover's moves and orange for the replies.}
\fi
\label{fig:encoding}
\end{figure}
\fi


\paragraph{Inputs of this task.}
The task is to predict whether a known tactic is found, so every model is
given the solution. The auxiliary planes mark the en passant square, the
opponent's mistake, the first three moves of the solver and the two replies
between them, each move by its origin and destination square, and one further
plane marks every square that a later move of the solution leaves or reaches.
A network therefore sees where the pieces move. It is not told whether a move
is a capture, a check or a sacrifice, and has to infer this from the board.
The hand-crafted features state these properties (Table~\ref{tab:hand}).

Besides the castling rights, the scalar context features are the two
ratings, the increment, the ply, the number of moves that the solver has to
find, the mover's win probability before and after the mistake, and three
times, each as the logarithm of one plus the number of seconds: the time
remaining on the mover's clock, the time remaining on the opponent's clock
and the initial time of the time control. The tables and figures call the
win probability after the mistake the win probability before the move. The
two win probabilities together tell a model how much the mistake changed the
evaluation. The move played and the time spent on it are known only
afterwards, and no model receives them. The puzzle's rating and its theme
tags are not inputs of any trained model either, because a tactic that arises
in a user's own game has neither. They appear in one baseline
(Section~\ref{sec:models}) and in the error analysis
(Section~\ref{sec:error}).

\subsection{Models and training}
\label{sec:models}

% Models and training: the same sequence of models for every task, so each
% model paper inputs this file under its own heading.
% The prose follows docs/writing.md.
Every task is given the same sequence of models, so that each step of the
sequence measures the contribution of one component. A model is fitted on the
training set, selected on the validation set and scored once on the test set.
Table~\ref{tab:runs} lists the models of this task.

\paragraph{Constant baseline.}
\run{constant} predicts the same value for every example: the mean of the
training targets (for a binary target, the base rate), or their median when
the primary metric is an absolute error. A useful model has to improve on it.

\paragraph{Gradient-boosted decision trees.}
\ifcontextgbdt
\run{gbm-context} is a gradient-boosted decision tree model (GBDT), trained
with LightGBM~\cite{lightgbm} on the scalar context features alone. It never
sees the board, and so measures how much of the target the context already
determines. \run{gbm-hand} adds the hand-crafted features.
\else
\run{gbm-hand} is a gradient-boosted decision tree model (GBDT), trained with
LightGBM~\cite{lightgbm} on the scalar context features and the hand-crafted
features.
\fi
We refer to \run{gbm-hand} as \emph{the GBDT baseline}: every other model is
compared with it on the same test examples. GBDT is trained for at most
\V{method.gbm.rounds} boosting rounds with up to \V{method.gbm.leaves} leaves
per tree and a learning rate of \V{method.gbm.rate}. Each tree is fitted on a
random \V{method.gbm.rows.percent}\% of the examples and
\V{method.gbm.features.percent}\% of the features, and boosting stops when the
validation score has not improved for \V{method.gbm.patience} rounds.

\paragraph{Neural networks.}
A network consists of an encoder, which maps the square tokens, the auxiliary
planes and the context features to a vector, and a prediction head, a
two-layer perceptron that maps this vector to the output. We use two encoder
architectures, so that the results show how much the architecture matters for
the task.

The \emph{transformer} encoder~\cite{attention} treats each square as a token.
The embedding of a token is the sum of a learned embedding of its piece, a
learned embedding of its square and a linear projection of the square's
auxiliary planes. The context features are projected by a small perceptron
and enter as one additional token. Every token attends to every other token
in every layer, so distant squares interact from the first layer on. To give
attention the geometry of the board, each attention head adds a learned bias
that depends on the file and rank offset between two
squares~\cite{relative}: files, ranks, diagonals and knight moves all
correspond to fixed offsets. The layers use pre-layer normalisation, and the
encoder outputs the context token concatenated with the mean of the 64 square
tokens.

The \emph{residual convolutional network} (ResNet) is a stack of residual
blocks of $3\times3$ convolutions over the $8\times8$ board, the architecture
that AlphaZero introduced for chess~\cite{alphazero}. Its input is one binary
plane per token value together with the auxiliary planes. A linear projection
of the context features is added at every square after the first convolution,
and the encoder outputs the mean and the maximum over the squares of the
final feature map. Unlike attention, a convolution relates only neighbouring
squares, so the receptive field grows with depth.

\ifsequence
\paragraph{Sequence model.}
An example of this task is a sequence of positions, so a network has two
stages below its prediction head. One of the two encoders above is applied to
every position of the sequence, with the position's per-move features in
place of the context features, and produces one vector per position. A
two-layer transformer then operates on these vectors in the order in which
the moves were played: each vector is a token, a learned positional embedding
is added to it, and the game-level context features enter as one additional
token. The prediction head receives that token concatenated with the mean of
the position tokens. Sequences in a batch are padded to the length of the
longest, and padded elements are masked: they are not encoded and no token
attends to them, so a prediction does not depend on the other examples in
the batch.

Model sizes are hyperparameters. \run{conv-small} is small enough to train
on a CPU. An example consists of tens of positions, so
\run{transformer-small} and \run{transformer-medium} are trained on a GPU.
\run{transformer-medium-hand} is \run{transformer-medium} with the
hand-crafted features added to its context features.
Table~\ref{tab:settings} in the appendix lists the hyperparameters.
\else
Model sizes are hyperparameters. \run{conv-small} and \run{transformer-small}
are small enough to train on a CPU. \run{transformer-medium} is trained on a
GPU, and \run{transformer-medium-hand} is the same network with the
hand-crafted features added to its context features.
Table~\ref{tab:settings} in the appendix lists the hyperparameters.
\fi

\paragraph{Training.}
All networks are trained in the same way: with AdamW~\cite{adamw}, a learning
rate that increases linearly over the first \V{networks.warmup.percent}\% of
the optimisation steps and then follows a cosine schedule to zero, gradient
clipping, and mixed precision on a GPU. The loss matches the primary metric:
mean squared error for a target scored by RMSE, a Huber loss for a target
scored by MAE, and binary cross-entropy for a probability. A regression
target is standardised with the mean and standard deviation of the training
set. After every epoch the network is evaluated on the validation set. We
keep the weights of the best epoch, and stop training when
\V{networks.patience} consecutive epochs have not improved on it.
\ifmore
A model trained on a larger training set is given proportionally fewer
epochs, and a correspondingly shorter patience
(Section~\ref{sec:ablations}).
\fi

\begin{table}[htbp]\centering\footnotesize
\block{runs}
\caption{The models, best first. \emph{Selected} is the number of trees or
the training epoch chosen on the validation set. \emph{Time} is the total
wall-clock time for training and evaluation.}
\label{tab:runs}
\end{table}


\paragraph{Rating-based baseline.}
\run{puzzle-rating} is a baseline that uses only the mover's rating $r_m$ and
the puzzle's rating $r_p$. It applies the Elo formula for the expected score
to the two,
\[
p = \frac{1}{1 + 10^{-(r_m - r_p)/400}},
\]
and its output is calibrated like that of every other model
(Section~\ref{sec:protocol}). The slope and the intercept of the calibration
are its only fitted parameters. Because the log-odds of the formula are
proportional to the rating difference, the slope rescales this difference and
the intercept shifts it, which allows for the different scales of game
ratings and puzzle ratings. The puzzle's rating summarises the attempts of
many solvers and is the most direct measurement available of how hard the
tactic is as a puzzle. The baseline shows how much of the outcome in a game
the two ratings explain, and its slope shows how well the formula transfers
from puzzles to games.

\section{Evaluation Protocol}
\label{sec:protocol}

\paragraph{Metrics.}
The primary metric, fixed before the first full training run, is the log
loss on the test set: the negative natural logarithm of the probability that
a model assigned to the observed outcome, averaged over the test positions.
Lower is better. We chose it because the application would display the
prediction as a probability, and the log loss is a proper scoring rule that
penalises a probability both for failing to separate the outcomes and for
being miscalibrated. We report four secondary metrics: the Brier score, the
expected calibration error, which is the absolute difference between the
mean predicted probability and the observed share of tactics found, averaged
over ten bins of equal size formed by the predicted probability, and the
areas under the precision-recall curve (PR-AUC) and under the receiver
operating characteristic curve (ROC-AUC), which measure only how well a model
ranks the positions. We also report results separately for \V{bands.count}
bands of the mover's rating.

% Evaluation protocol: the same for every task. Each model paper names its
% own metrics under its section heading and inputs the rest from here.
% The prose follows docs/writing.md.
\paragraph{Confidence intervals.}
Examples from the same game are not independent, so we resample whole games.
The test games are drawn with replacement \V{method.resamples} times, the
metric is computed each time on the examples of the games drawn, and the
2.5th and 97.5th percentiles of these values form the 95\% confidence
interval (CI) that we report with every score.

\paragraph{Paired comparisons.}
To compare two models we compute both scores on the same resampled games and
take their difference. We call this difference the \emph{gain} of the first
model over the second, and sign it so that a positive gain means that the
first model is better. A gain is \emph{statistically significant} when its
95\% CI excludes zero, and is then set in bold in the tables. Unless stated
otherwise, the second model is the GBDT baseline. Two models tend to do well
and badly on the same games, so a paired gain can be significant even when
the two models' separate intervals overlap. These intervals reflect the
sampling of test games only. Each network was trained with a single random
seed, and the variation between training runs is not included.

\paragraph{Model selection.}
\ifchance
The number of trees, the training epoch whose weights are kept and the
calibration parameters are chosen on the validation set.
\else
The number of trees and the training epoch whose weights are kept are chosen
on the validation set.
\fi
The test set is used once per model, for the scores reported here.

\ifchance
\paragraph{Calibration.}
Before a model is scored, its predicted probabilities are rescaled by Platt
scaling~\cite{platt}: a logistic function of the log-odds whose slope and
intercept are fitted on the validation set. The transformation is monotone,
so it leaves a model's ranking of the examples unchanged and affects only its
log loss and its calibration. What remains of a model's miscalibration after
this step cannot be removed by a linear correction of its log-odds.
\fi

\ifpairs
\paragraph{Acceptance criterion.}
Section~\ref{sec:ablations} evaluates modifications of the models. A
modification is \emph{accepted} only if the modified model is better than the
unmodified one on the validation set and its paired gain on the test set is
statistically significant. A modification that is not accepted is still
reported, as a measurement of its effect.
\fi

\paragraph{Tables and figures.}
A model has the same colour and marker in every figure, and the GBDT baseline
is always drawn in orange. The values plotted in a figure are also given in a
table, in the text or in the appendix. All numbers, tables and figures are
generated from the committed result files by the command that builds the
paper.


\section{Results}
\label{sec:results}

\subsection{Main results}

Table~\ref{tab:leaderboard} and Figure~\ref{fig:scores} give the main
results. The constant baseline predicts the share of tactics found in the
training set, \V{floor.constant}, for every position, and has a log loss of
\V{floor.score}. GBDT on the context features alone reaches
\V{run.gbm-context.score} and thereby achieves \V{context.lift.share}\% of
the improvement of the GBDT baseline over the constant baseline. The GBDT
baseline, with the hand-crafted features of the position and of the
solution, reaches \V{reference.score}. The larger part of what is predicted
thus comes from the description of the tactic, and not from the ratings, the
clocks and the evaluations.

\begin{table}[htbp]\centering\small
\block{leaderboard}
\caption{Main results on the \V{data.rows.test} test positions, best first.
The gain is the paired difference in log loss from \V{task.reference} on the
same resampled games. It is positive when the model is better, and bold when
statistically significant.}
\label{tab:leaderboard}
\end{table}

\begin{figure}[htbp]\centering
\block{scores}
\caption{Test log loss of every model with its 95\% CI. The vertical line
marks the GBDT baseline.}
\label{fig:scores}
\end{figure}

Networks trained from random initialisation on the board, the marked squares
of the solution and the context features do not reach the GBDT baseline.
Compared with it, \V{ladder.behind.of} of the fixed sequence are
significantly worse (Figure~\ref{fig:gains}): the small ResNet (gain
\V{run.conv-small.gain}; 95\% CI \V{run.conv-small.gain.low} to
\V{run.conv-small.gain.high}), the small transformer
(\V{run.transformer-small.gain}; \V{run.transformer-small.gain.low} to
\V{run.transformer-small.gain.high}) and the medium transformer
(\V{run.transformer-medium.gain}; \V{run.transformer-medium.gain.low} to
\V{run.transformer-medium.gain.high}). The three models have almost the same
log loss (\V{run.conv-small.score}, \V{run.transformer-small.score} and
\V{run.transformer-medium.score}), although the medium transformer has
\V{run.transformer-medium.parameters.millions} million parameters and the
small models have \V{run.conv-small.parameters.millions} and
\V{run.transformer-small.parameters.millions} million. All three are far
better than GBDT on the context features alone, which they also receive. From
the board and the marked squares, a network thus learns most of what the
hand-crafted features give GBDT, but not all of it, and neither the size of
the network nor the choice between a convolutional and a transformer encoder
changes this.

\begin{figure}[htbp]\centering
\block{gains}
\caption{Paired gain in log loss of each network over the GBDT baseline, with
its 95\% CI. A model whose interval lies entirely to the right of the line is
significantly better than the baseline, and one whose interval lies entirely
to the left is significantly worse.}
\label{fig:gains}
\end{figure}

The hand-crafted features change this result. \run{transformer-medium-hand},
the medium transformer with the hand-crafted features added to its context
features, is significantly better than the GBDT baseline (gain
\V{run.transformer-medium-hand.gain}; 95\% CI
\V{run.transformer-medium-hand.gain.low} to
\V{run.transformer-medium-hand.gain.high}).
\ifpairs
Compared with \run{transformer-medium}, it is
\V{pair.transformer-medium-hand.transformer-medium.verdict} on the test set
(gain \V{pair.transformer-medium-hand.transformer-medium.gain}; 95\% CI
\V{pair.transformer-medium-hand.transformer-medium.gain.low} to
\V{pair.transformer-medium-hand.transformer-medium.gain.high}).
\fi
A network that reads the board therefore adds information to the hand-crafted
features, but on this training set it does not learn from the board all that
the features provide.

The remaining models in Table~\ref{tab:leaderboard} are the rating-based
baseline, which the next subsection discusses, and modifications that passed
the acceptance criterion of Section~\ref{sec:protocol}, which
Section~\ref{sec:ablations} describes\ifchangestart: two transformers with a
pretrained encoder, which are the two best models\fi\ifchangeconv, a
ResNet with hand-crafted features\fi\ifchangeepochs, a transformer with a
longer schedule\fi\ifmore, and models trained on a larger training
set\fi. Compared with the GBDT baseline, \V{networks.ahead.of} on the
leaderboard are significantly better. The best model's gain is
\V{best.gain.percent}\% of the GBDT baseline's log loss, and
\V{best.gain.of-lift}\% of the improvement of the GBDT baseline over the
constant baseline. GBDT on the context and the hand-crafted features
therefore accounts for most of what is predicted, and the best network adds a
smaller improvement.

\subsection{The two ratings as a predictor}
\label{sec:rating-baseline}

The rating-based baseline \run{puzzle-rating} has a log loss of
\V{run.puzzle-rating.score}. It is significantly worse than the GBDT baseline
(gain \V{run.puzzle-rating.gain}; 95\% CI \V{run.puzzle-rating.gain.low} to
\V{run.puzzle-rating.gain.high}), which is not given the puzzle's rating. Its
log loss is close to that of GBDT on the context features alone
(\V{run.gbm-context.score}), and the intervals of the two models' gains
overlap (Table~\ref{tab:leaderboard}). The two metrics that depend only on
the ranking of the positions do not agree on the order of the two models: the
rating-based baseline has the lower PR-AUC (\V{run.puzzle-rating.pr-auc},
compared with \V{run.gbm-context.pr-auc}) and the higher ROC-AUC
(\V{run.puzzle-rating.roc-auc}, compared with \V{run.gbm-context.roc-auc};
Table~\ref{tab:metrics}). The puzzle's rating and the mover's rating together thus predict the outcome in a
game about as well as the context of the game does without any knowledge of
the tactic beyond the length of its solution, and much less well than a
description of the tactic.

The calibration step shows how far the Elo formula is from describing games.
The fitted slope is \V{run.puzzle-rating.platt.slope}. A slope of one would
mean that a difference between the two ratings has the effect that the
formula assigns to it, up to a shift between the two rating scales. The
fitted value means that the share of tactics found depends on the difference
between the two ratings much less than the formula implies. Even after this
correction the baseline is the least well calibrated of all models, with an
expected calibration error of \V{run.puzzle-rating.calibration-error}
(Table~\ref{tab:metrics}).\iflook{} The observed shares show the weak
dependence directly (Figure~\ref{fig:grid-ratings} in
Section~\ref{sec:error}). Movers rated below 1700 found
\V{grid.ratings.-1700.2200+.happened} of the tactics whose puzzles are rated
2200 and above. Such a mover is rated at least 500 points below the puzzle,
for which the uncalibrated formula gives an expected score of at most 5.3\%.
Movers rated 2300 and above found \V{grid.ratings.2300+.-1000.happened} of
the tactics whose puzzles are rated below 1000. Such a mover is rated more
than 1300 points above the puzzle, for which the formula gives more than
99.9\%.\fi

This result does not mean that the puzzle's rating is uninformative. The
share of tactics found falls steadily as the puzzle's rating rises, and the
trained models, which do not receive the rating, are miscalibrated with
respect to it (Section~\ref{sec:error}).

\subsection{Secondary metrics}

Table~\ref{tab:metrics} gives all metrics for every model. The Brier score of
the best model is \V{best.brier}, compared with \V{reference.brier} for the
GBDT baseline and \V{floor.brier} for the constant baseline. Its PR-AUC is
\V{best.pr-auc}, compared with \V{reference.pr-auc} and \V{floor.pr-auc}.
The PR-AUC of a model that ranks at random equals the share of tactics found
in the test set. Because this share is high, all models lie in a narrow
range above it. The ROC-AUC of the best model is \V{best.roc-auc}, compared
with \V{reference.roc-auc} for the GBDT baseline. A value this far from one
means that the outcome of a single position remains uncertain: the models
separate tactics that are usually found from tactics that are often missed,
and they do not predict what an individual player does.

\begin{table}[htbp]\centering\small
\block{metrics}
\caption{All metrics for every model on the test set. For the log loss, the
Brier score and the calibration error lower is better. For PR-AUC and ROC-AUC
higher is better.}
\label{tab:metrics}
\end{table}

\subsection{Calibration}

A predicted probability is useful as a number only if it is calibrated: among
the positions for which a model predicts a probability of 70\%, about seven
tactics in ten should be found. In Figure~\ref{fig:calibration} the test
positions are divided into \V{best.calibration.groups} bins of equal size by
a model's predicted probability. For the best model, the mean predicted
probability is \V{best.calibration.lowest}\% in the lowest bin and
\V{best.calibration.highest}\% in the highest, and in no bin does it differ
from the observed share of tactics found by more than
\V{best.calibration.widest-gap} percentage points. For the GBDT baseline the
predictions range from \V{reference.calibration.lowest}\% to
\V{reference.calibration.highest}\%, and the largest difference is
\V{reference.calibration.widest-gap} points. Table~\ref{tab:calibration} in
the appendix gives the values for every bin.

\begin{figure}[htbp]\centering
\block{calibration}
\caption{Reliability diagram. The test positions are divided into
\V{best.calibration.groups} bins of equal size by a model's predicted
probability, and each point shows the observed share of tactics found in a
bin against its mean predicted probability. Points on the diagonal are
perfectly calibrated.}
\label{fig:calibration}
\end{figure}

The calibration step of Section~\ref{sec:protocol} changes the networks more
than GBDT. The fitted slope is \V{best.platt.slope} for the best model and
\V{run.transformer-medium-hand.platt.slope} for
\run{transformer-medium-hand}. A slope below one means that the uncalibrated
predictions were too extreme, and the calibration step makes them less so.
For the GBDT baseline the slope is \V{reference.platt.slope} and the
intercept \V{reference.platt.intercept}, so this model was close to
calibrated without correction. The expected
calibration error after calibration is \V{best.calibration-error} for the
best model and \V{reference.calibration-error} for the GBDT baseline.
Calibration over the whole test set does not imply calibration within every
subgroup, and Section~\ref{sec:error} returns to this point.

\iflook
\subsection{The two ends of the ranking}

The application would use the prediction at both ends of its range. A missed
tactic with a high predicted probability is a lapse, and a missed tactic with
a low predicted probability is one that the player could not be expected to
find. Table~\ref{tab:ranking-top} gives the share of tactics found among the
test positions with the highest predictions of the best model, and
Table~\ref{tab:ranking-bottom} the share missed among those with the lowest.

Among the 10\% of test positions with the highest predictions,
\V{best.top.10.rate}\% of the tactics were found, and among the top 1\%,
\V{best.top.1.rate}\%. Because \V{label.rate.test}\% of all tactics in the
test set are found, the lift at the top of the ranking is necessarily small:
it is \V{best.top.10.lift} for the top 10\%. A miss in these
positions is nevertheless a rare event, which is what the application needs
in order to call it a lapse.

\begin{table}[htbp]\centering\small
\block{ranking-top}
\caption{The top of the ranking. For a share of the test positions taken from
the highest predicted probabilities: \emph{precision} is the share of these
positions in which the tactic was found, \emph{recall} is the share of all
found tactics that they contain, and \emph{lift} is the precision divided by
the share found in the whole test set.}
\label{tab:ranking-top}
\end{table}

\begin{table}[htbp]\centering\small
\block{ranking-bottom}
\caption{The bottom of the ranking. For a share of the test positions taken
from the lowest predicted probabilities: \emph{precision} is the share of
these positions in which the tactic was missed, \emph{recall} is the share of
all missed tactics that they contain, and \emph{lift} is the precision
divided by the share missed in the whole test set.}
\label{tab:ranking-bottom}
\end{table}

The bottom of the ranking is more informative. Among the 10\% of test
positions with the lowest predictions, \V{best.bottom.10.rate}\% of the
tactics were missed, which is \V{best.bottom.10.lift} times the share missed
in the whole test set, and these positions contain
\V{best.bottom.10.recall}\% of all missed tactics. Among the bottom 1\%,
\V{best.bottom.1.rate}\% were missed. The remaining tactics in these
positions were found, so a low prediction marks a tactic that is usually
missed and not one that is never found. The GBDT baseline is slightly less
accurate at this end: among its own bottom 10\%,
\V{reference.bottom.10.rate}\% of the tactics were missed.
\fi

\subsection{Results by rating band}

Table~\ref{tab:bands} gives the log loss by band of the mover's rating. The
log loss of the constant baseline falls from \V{band.1.floor} in the band
\V{band.1.name} to \V{band.4.floor} in the band \V{band.4.name}.\iflook{} The reason is that the share of tactics found
rises with the rating, from \V{slice.band.-1700.happened} to
\V{slice.band.2100-2299.happened} in these two bands, and a constant
prediction above one half is scored better where more tactics are found.\fi{}
The log loss of every model depends on this base rate, so scores should be
compared between models within a band, and not between bands.

\begin{table}[htbp]\centering\small
\block{bands-table}
\caption{Test log loss by band of the mover's rating. The share of tactics
found differs between bands, so values are comparable within a column and
not between columns.}
\label{tab:bands}
\end{table}

The best model has a lower log loss than the GBDT baseline in
\V{bands.best-ahead.of}. The difference is largest in the band
\V{bands.lead.widest.name} (\V{bands.lead.widest}) and smallest in the band
\V{bands.lead.narrowest.name} (\V{bands.lead.narrowest};
Figure~\ref{fig:bands}).\iflook{} With a paired comparison within each band,
the gain is statistically significant in \V{slice.band.ahead.of}
(Table~\ref{tab:slice-band} in the appendix).\fi

\begin{figure}[htbp]\centering
\block{bands}
\caption{Difference in log loss between each network and the GBDT baseline,
by band of the mover's rating (the differences of the values in
Table~\ref{tab:bands}, without intervals). Above the dashed line the network
is better.}
\label{fig:bands}
\end{figure}

\subsection{Training curves}

Figure~\ref{fig:curves} shows the validation log loss of each network of the
fixed sequence after every epoch. By this measure
\V{ladder.at-last-epoch.of} were still improving when training ended:
\run{transformer-small} had its best validation log loss at epoch
\V{run.transformer-small.best-epoch} of \V{run.transformer-small.epochs} and
\run{transformer-medium} at epoch \V{run.transformer-medium.best-epoch} of
\V{run.transformer-medium.epochs}. The other two had passed their best epoch
when training ended: \run{conv-small} at \V{run.conv-small.best-epoch} of
\V{run.conv-small.epochs} and \run{transformer-medium-hand} at
\V{run.transformer-medium-hand.best-epoch} of
\V{run.transformer-medium-hand.epochs}.
\ifchangeepochs
Section~\ref{sec:ablations} shows that a longer schedule gives the medium
transformer without hand-crafted features a small significant gain, and that
it does not improve the one with hand-crafted features.
\fi

\begin{figure}[htbp]\centering
\block{curves}
\caption{Validation log loss of each network of the fixed sequence after
every epoch. The dashed line is the GBDT baseline's validation log loss.}
\label{fig:curves}
\end{figure}

\subsection{Feature importance of the GBDT baseline}

The GBDT baseline consists of \V{reference.trees} trees over
\V{reference.features} features. Figure~\ref{fig:leans} shows the ten
features with the largest shares of the total split gain. The importance is spread over many features,
and none has more than \V{reference.leans.1.share}\% of the gain. The two
most important are the \V{reference.leans.1.feature}
(\V{reference.leans.1.share}\%) and whether the \V{reference.leans.2.feature}
(\V{reference.leans.2.share}\%), followed by the \V{reference.leans.3.feature}
(\V{reference.leans.3.share}\%), the \V{reference.leans.4.feature}
(\V{reference.leans.4.share}\%) and the \V{reference.leans.5.feature}
(\V{reference.leans.5.share}\%). Split gain shows which features the model
used. It does not identify what causes a tactic to be missed, because
correlated features divide the credit between them arbitrarily.

\begin{figure}[htbp]\centering
\block{leans}
\caption{The ten most important features of the GBDT baseline, by share of
the total split gain.}
\label{fig:leans}
\end{figure}

\ifchanges
\section{Ablations}
\label{sec:ablations}

% Ablations: what a modification is, the table of paired comparisons and the
% two figures of validation curves. The same for every task, so a model paper
% inputs this file at the start of its ablation section and then reports its
% own findings. Each part appears once the task has the runs it needs, and the
% list below needs at least one of its items.
% The prose follows docs/writing.md.
\ifpairs
The main results compare complete models. An ablation isolates one factor: it
compares a model with its \emph{base}, a model that differs from it in a
single respect. Table~\ref{tab:pairs} lists every such comparison for this
task, with the gain of the modified model over its base on the validation set
and its paired gain on the test set, computed on the same resampled test
games as in the main results. The acceptance criterion of
Section~\ref{sec:protocol} is applied to these pairs. The modifications are
the following.
\begin{itemize}
\ifchangehand
\item \textbf{Hand-crafted features} (suffix \code{-hand}): the network
receives the hand-crafted features as additional context features.
\fi
\ifchangeepochs
\item \textbf{Longer schedule} (suffix \code{-long}): twice the training
epochs, with the same warm-up fraction and cosine schedule.
\fi
\ifchangestart
\item \textbf{Pretraining} (suffix \code{-pretrained}): the transformer
encoder is initialised from a transformer of the same size that was trained
to predict the engine's win probability for positions from the Lichess
evaluation database~\cite{lichessdb}. The piece and square embeddings and all
transformer layers are copied. The prediction head is initialised randomly,
as are the projections of the auxiliary planes that the pretraining task does
not have and, unless the two tasks have the same context features, the
projection of the context features. Training then follows the standard
schedule.
\fi
\ifchangeattacks
\item \textbf{Attack-aware attention bias} (suffix \code{-attacks}): the
relative position bias encodes where two squares lie and not what stands
between them. With this modification each attention head learns four
additional biases for a pair of squares, applied when the piece on one square
attacks the other, when it is attacked from the other, and for the same two
relations through exactly one intervening piece, as in a pin or a skewer. The
relations are computed from the tokens without search, and the biases are
initialised to zero.
\fi
\ifchangegeometry
\item \textbf{No relative position bias} (suffix \code{-no-geometry}): the
relative position bias is removed. This configuration measures the
contribution of the bias and is not a candidate model.
\fi
\ifchangeconv
\item \textbf{ResNet encoder} (prefix \code{conv-medium}): a ResNet with
squeeze-and-excitation in place of the transformer, trained at the higher
peak learning rate used for the convolutional models.
\fi
\ifchangehybrid
\item \textbf{Hybrid encoder} (prefix \code{hybrid-medium}): a few residual
convolutional blocks followed by transformer layers over the resulting
square representations.
\fi
\ifchangelarger
\item \textbf{Larger transformer} (prefix \code{transformer-large}): more layers and
a larger width, with the training schedule unchanged.
\fi
\ifgames
\item \textbf{Training-set size}: the same model trained on a fraction or a
multiple of the training games (Section~\ref{sec:training-size}).
\fi
\iffeatures
\item \textbf{Feature groups}: a model trained without one group of
hand-crafted features at a time (Section~\ref{sec:feature-groups}).
\fi
\end{itemize}

\longblock{pairs}{Paired comparisons. Each row compares a model with a
baseline, in most rows the model it modifies. \emph{Change} states how the
two differ. The gain is that of the model over the baseline and is positive
when the model is better. On
the validation set it is the difference between the two scores. On the test
set it is the paired gain, with its 95\% CI, and is set in bold when
statistically significant.}{tab:pairs}
\fi

\ifvariants
Figure~\ref{fig:variants} shows the validation curves of each model of the
fixed sequence whose training schedule or initialisation was modified,
together with the modified models. A modified model has the colour of its
base and a broken line: long dashes for a longer schedule and dots for a
pretrained encoder.

\begin{figure}[htbp]\centering
\block{variants}
\caption{Validation curves of a model of the fixed sequence (solid line) and
of the models that modify its training schedule or initialisation (broken
lines), after every epoch. One panel per base model, on a common scale. The
dashed horizontal line is the GBDT baseline's validation score.}
\label{fig:variants}
\end{figure}
\fi

\ifshapes
Figure~\ref{fig:shapes} shows the validation curves of the models that
change the architecture and nothing about the training data, each in its own
panel beside its base.

\begin{figure}[htbp]\centering
\block{shapes}
\caption{Validation curves of a model with a modified architecture (broken
line) and of its base (solid line), after every epoch. One panel per modified
model, on a common scale. The dashed horizontal line is the GBDT baseline's
validation score.}
\label{fig:shapes}
\end{figure}
\fi


\ifchangehand
\subsection{Hand-crafted features}

Among the networks trained from random initialisation, only those that
receive the hand-crafted features are better than the GBDT baseline. Compared
with \run{transformer-medium}, \run{transformer-medium-hand} is
\V{pair.transformer-medium-hand.transformer-medium.valid.verdict} than its
base on the validation set and
\V{pair.transformer-medium-hand.transformer-medium.verdict} on the test set
(gain \V{pair.transformer-medium-hand.transformer-medium.gain} log loss;
95\% CI \V{pair.transformer-medium-hand.transformer-medium.gain.low} to
\V{pair.transformer-medium-hand.transformer-medium.gain.high}), so the
modification is \V{pair.transformer-medium-hand.transformer-medium.kept}. The
features move the medium transformer from significantly below the GBDT
baseline to significantly above it: its gains over the baseline are
\V{run.transformer-medium.gain} without them and
\V{run.transformer-medium-hand.gain} with them.
\iffeatures
Section~\ref{sec:feature-groups} shows which of the features matter most to
GBDT.
\fi
\fi

\ifchangeepochs
\subsection{Training schedule}

Doubling the number of epochs, from \V{run.transformer-medium.epochs} to
\V{run.transformer-medium-long.epochs}, helps the medium transformer without
hand-crafted features and not the one with them. Without hand-crafted
features, \run{transformer-medium-long} is
\V{pair.transformer-medium-long.transformer-medium.valid.verdict} than its
base on the validation set and
\V{pair.transformer-medium-long.transformer-medium.verdict} on the test set
(gain \V{pair.transformer-medium-long.transformer-medium.gain}; 95\% CI
\V{pair.transformer-medium-long.transformer-medium.gain.low} to
\V{pair.transformer-medium-long.transformer-medium.gain.high}), so the
modification is \V{pair.transformer-medium-long.transformer-medium.kept}. The
gain is small beside that from the hand-crafted features
(\V{pair.transformer-medium-hand.transformer-medium.gain}), and the model
remains significantly worse than the GBDT baseline (gain
\V{run.transformer-medium-long.gain}; 95\% CI
\V{run.transformer-medium-long.gain.low} to
\V{run.transformer-medium-long.gain.high}). With hand-crafted features,
\run{transformer-medium-hand-long} is
\V{pair.transformer-medium-hand-long.transformer-medium-hand.valid.verdict}
than its base on the validation set and
\V{pair.transformer-medium-hand-long.transformer-medium-hand.verdict} on the
test set (gain
\V{pair.transformer-medium-hand-long.transformer-medium-hand.gain}; 95\% CI
\V{pair.transformer-medium-hand-long.transformer-medium-hand.gain.low} to
\V{pair.transformer-medium-hand-long.transformer-medium-hand.gain.high}),
and the modification is
\V{pair.transformer-medium-hand-long.transformer-medium-hand.kept}.

The training curves agree with this difference
(Figure~\ref{fig:variants}). Under the standard schedule,
\run{transformer-medium} had its best validation log loss in the final
epoch, \V{run.transformer-medium.best-epoch} of
\V{run.transformer-medium.epochs}, whereas \run{transformer-medium-hand} had
it at epoch \V{run.transformer-medium-hand.best-epoch} of
\V{run.transformer-medium-hand.epochs}. Under the longer schedule neither
model used all its epochs: \run{transformer-medium-long} had its best
validation log loss at epoch \V{run.transformer-medium-long.best-epoch} of
\V{run.transformer-medium-long.epochs} and
\run{transformer-medium-hand-long} at epoch
\V{run.transformer-medium-hand-long.best-epoch}, and each was stopped
\V{networks.patience} epochs later by early stopping.
\fi

\ifchangestart
\subsection{Pretraining}

A pretrained encoder improves both medium transformers. Without hand-crafted
features, \run{transformer-medium-pretrained} is
\V{pair.transformer-medium-pretrained.transformer-medium.valid.verdict} than
its base on the validation set and
\V{pair.transformer-medium-pretrained.transformer-medium.verdict} on the
test set (gain \V{pair.transformer-medium-pretrained.transformer-medium.gain}
log loss; 95\% CI
\V{pair.transformer-medium-pretrained.transformer-medium.gain.low} to
\V{pair.transformer-medium-pretrained.transformer-medium.gain.high}), and
the modification is
\V{pair.transformer-medium-pretrained.transformer-medium.kept}. With
hand-crafted features, \run{transformer-medium-hand-pretrained} is
\V{pair.transformer-medium-hand-pretrained.transformer-medium-hand.valid.verdict}
than its base on the validation set and
\V{pair.transformer-medium-hand-pretrained.transformer-medium-hand.verdict}
on the test set (gain
\V{pair.transformer-medium-hand-pretrained.transformer-medium-hand.gain};
95\% CI
\V{pair.transformer-medium-hand-pretrained.transformer-medium-hand.gain.low}
to
\V{pair.transformer-medium-hand-pretrained.transformer-medium-hand.gain.high}),
and the modification is
\V{pair.transformer-medium-hand-pretrained.transformer-medium-hand.kept}.
The best model of this task is \V{best.name}. The paired gain without
hand-crafted features is the largest of any modification. It exceeds the
gain from the hand-crafted features themselves
(\V{pair.transformer-medium-hand.transformer-medium.gain}), and the intervals
of the two gains do not overlap. The paired gain with hand-crafted features
is the largest among the modifications of \run{transformer-medium-hand}
(Table~\ref{tab:pairs}).

Pretraining changes the comparison with the GBDT baseline. Trained from
random initialisation, the transformer without hand-crafted features is
significantly worse than the baseline (gain
\V{run.transformer-medium.gain}). With a pretrained encoder it is
significantly better (gain \V{run.transformer-medium-pretrained.gain}; 95\%
CI \V{run.transformer-medium-pretrained.gain.low} to
\V{run.transformer-medium-pretrained.gain.high}).
\ifchangeepochs
Apart from the hand-crafted features, no other modification of this model
has that effect. With the longer
schedule its gain over the baseline is \V{run.transformer-medium-long.gain}
(95\% CI \V{run.transformer-medium-long.gain.low} to
\V{run.transformer-medium-long.gain.high})\ifchangeattacks, and with the
attack-aware attention bias it is \V{run.transformer-medium-attacks.gain}
(\V{run.transformer-medium-attacks.gain.low} to
\V{run.transformer-medium-attacks.gain.high})\fi.
\fi

With a pretrained encoder the hand-crafted features contribute less, but
they still contribute. The pretrained model with hand-crafted features has a
test log loss of \V{run.transformer-medium-hand-pretrained.score} and the
pretrained model without them one of
\V{run.transformer-medium-pretrained.score}, a difference of
\V{run.transformer-medium-hand-pretrained.over-twin}, where the paired gain
from the features is \V{pair.transformer-medium-hand.transformer-medium.gain}
from random initialisation. The two pretrained models were not compared
directly on resampled games, but the intervals of their gains over the GBDT
baseline do not overlap (Table~\ref{tab:leaderboard}). A plausible
explanation is that the pretraining task evaluates a position and involves no
line of moves, so pretraining teaches the encoder nothing about a given
solution, which is what many of the hand-crafted features describe.

The pretrained model with hand-crafted features had its best validation log
loss at epoch \V{run.transformer-medium-hand-pretrained.best-epoch} of
\V{run.transformer-medium-hand-pretrained.epochs}, and the pretrained model
without them at epoch \V{run.transformer-medium-pretrained.best-epoch} of
\V{run.transformer-medium-pretrained.epochs} (Figure~\ref{fig:variants}).
Both were stopped \V{networks.patience} epochs later by early stopping.
\fi

\ifshapes
\subsection{Architecture}

The architecture ablations keep the training data, the prediction head and
the number of epochs of a medium transformer and change the encoder. All of
these models are trained from random initialisation, as their bases are.
\begin{itemize}
\ifchangeconv
\item \textbf{ResNet.} \run{conv-medium-hand} has
\V{run.conv-medium-hand.parameters.millions} million parameters, compared
with \V{run.transformer-medium-hand.parameters.millions} million for the
transformer, and is trained with a peak learning rate of
\V{run.conv-medium-hand.lr} instead of \V{run.transformer-medium-hand.lr}.
Compared with \run{transformer-medium-hand} it is
\V{pair.conv-medium-hand.transformer-medium-hand.valid.verdict} on the
validation set and
\V{pair.conv-medium-hand.transformer-medium-hand.verdict} on the test set
(gain \V{pair.conv-medium-hand.transformer-medium-hand.gain}; 95\% CI
\V{pair.conv-medium-hand.transformer-medium-hand.gain.low} to
\V{pair.conv-medium-hand.transformer-medium-hand.gain.high}), and the
modification is \V{pair.conv-medium-hand.transformer-medium-hand.kept}.
\fi
\ifchangehybrid
\item \textbf{Hybrid.} \run{hybrid-medium-hand} has
\V{run.hybrid-medium-hand.convs} residual convolutional blocks followed by
\V{run.hybrid-medium-hand.depth} transformer layers, for a total of
\V{run.hybrid-medium-hand.parameters.millions} million parameters. Compared
with \run{transformer-medium-hand} it is
\V{pair.hybrid-medium-hand.transformer-medium-hand.verdict} on the test set
(gain \V{pair.hybrid-medium-hand.transformer-medium-hand.gain}; 95\% CI
\V{pair.hybrid-medium-hand.transformer-medium-hand.gain.low} to
\V{pair.hybrid-medium-hand.transformer-medium-hand.gain.high}), and the
modification is \V{pair.hybrid-medium-hand.transformer-medium-hand.kept}.
\fi
\ifchangeattacks
\item \textbf{Attack-aware attention bias.} Compared with
\run{transformer-medium-hand}, \run{transformer-medium-hand-attacks} is
\V{pair.transformer-medium-hand-attacks.transformer-medium-hand.verdict} on
the test set (gain
\V{pair.transformer-medium-hand-attacks.transformer-medium-hand.gain}; 95\%
CI \V{pair.transformer-medium-hand-attacks.transformer-medium-hand.gain.low}
to
\V{pair.transformer-medium-hand-attacks.transformer-medium-hand.gain.high}),
and the modification is
\V{pair.transformer-medium-hand-attacks.transformer-medium-hand.kept}.
Compared with \run{transformer-medium}, which has no hand-crafted features,
\run{transformer-medium-attacks} is
\V{pair.transformer-medium-attacks.transformer-medium.verdict} on the test
set (gain \V{pair.transformer-medium-attacks.transformer-medium.gain}; 95\% CI
\V{pair.transformer-medium-attacks.transformer-medium.gain.low} to
\V{pair.transformer-medium-attacks.transformer-medium.gain.high}), and the
modification is
\V{pair.transformer-medium-attacks.transformer-medium.kept}.
\fi
\ifchangegeometry
\item \textbf{No relative position bias.} Compared with
\run{transformer-medium-hand},
\run{transformer-medium-hand-no-geometry} is
\V{pair.transformer-medium-hand-no-geometry.transformer-medium-hand.verdict}
on the test set (gain
\V{pair.transformer-medium-hand-no-geometry.transformer-medium-hand.gain};
95\% CI
\V{pair.transformer-medium-hand-no-geometry.transformer-medium-hand.gain.low}
to
\V{pair.transformer-medium-hand-no-geometry.transformer-medium-hand.gain.high}).
\fi
\end{itemize}

\ifchangeconv
One change to the architecture is accepted, the ResNet in place of the
transformer encoder. With a test log loss of \V{run.conv-medium-hand.score},
\run{conv-medium-hand} is the best model trained from random
initialisation\ifmore, including the transformer trained on twice the games
(\V{run.transformer-medium-hand-double.score})\fi.\ifchangestart{} Its paired
gain over the transformer is smaller than that of pretraining
(\V{pair.transformer-medium-hand-pretrained.transformer-medium-hand.gain}),
and the intervals of the two gains do not overlap (Table~\ref{tab:pairs}). A
ResNet with a pretrained encoder has not been trained.\fi
\fi
\ifchangehybrid{}
The hybrid encoder is significantly worse than the transformer, although it
remains significantly better than the GBDT baseline (gain
\V{run.hybrid-medium-hand.gain}; 95\% CI \V{run.hybrid-medium-hand.gain.low}
to \V{run.hybrid-medium-hand.gain.high}).
\fi
\ifchangeattacks
The attack-aware bias has no measurable effect on this task, with or without
the hand-crafted features\ifchangegeometry, and neither has the removal of
the relative position bias\fi. The paired gains of these models are close to
zero (Table~\ref{tab:pairs}). All networks of this task are trained with the
same random seed and both biases are initialised at zero, so each of these models starts
from the same function as its base. With the hand-crafted features, the
validation curves of the two variants can hardly be told from that of their
base (Figure~\ref{fig:shapes}), and the intervals of their paired gains are
very narrow. What the biases learn on this task therefore changes the
predictions very little. Nor does the attack-aware bias replace the hand-crafted
features: with it, the transformer without hand-crafted features has a test
log loss of \V{run.transformer-medium-attacks.score}, compared with
\V{run.transformer-medium-hand.score} for the transformer with the features.
\fi
Each of these models was trained once.
\fi

% Training-set size: the same ablation for every task, so a model paper
% inputs this file inside its ablation section and then reports its own
% findings. It appears once the task has the runs.
% The prose follows docs/writing.md.
\ifgames
\subsection{Training-set size}
\label{sec:training-size}

To estimate how much a model's accuracy depends on the amount of training
data, we retrain it on a half and on a quarter of the training games and
evaluate it on the full test set (Table~\ref{tab:games},
Figure~\ref{fig:games}). The subsets are selected by a hash of the game
identifier and are nested, so the quarter is contained in the half. A network
is given twice the epochs on half the games and four times the epochs on a
quarter, so the number of optimisation steps is constant and only the amount
of data varies.
\ifmore
We also train on more data. A larger sample of the same source is split in
the same way as the standard dataset, so the validation and test sets are
unchanged and no test game appears in the larger training set. A network
trained on a larger set is given proportionally fewer epochs, so that it
again takes the same number of optimisation steps as its base. Unlike a model
trained on a subset, a model trained on more data is a candidate model, and
the acceptance criterion applies to it.
\fi

\begin{table}[htbp]\centering\footnotesize
\block{study-games}
\caption{Training-set size. Each model is trained on a fraction of the
training games\ifmore, or on a multiple of them from the larger sample,\fi{}
and evaluated on the full test set. The gain is paired with the same model
trained on the standard training set. A negative gain means that the model is
worse, and bold marks statistical significance.}
\label{tab:games}
\end{table}

\begin{figure}[htbp]\centering
\block{study-games-figure}
\caption{Test score as a function of the number of training games. Each step
to the right doubles the training set.}
\label{fig:games}
\end{figure}
\fi


\ifgames
\paragraph{Smaller training sets.}
Both model classes are less accurate with less data, and the transformer
loses more. Relative to the same model trained on the full training set,
\run{transformer-medium-hand} has a paired gain of
\V{pair.transformer-medium-hand-half.transformer-medium-hand.gain} log loss
on half of the training games and of
\V{pair.transformer-medium-hand-quarter.transformer-medium-hand.gain} on a
quarter, and it is significantly worse in
\V{study.games.transformer-medium-hand.worse.of}. For GBDT the gains are
\V{pair.gbm-hand-half.gbm-hand.gain}
(\V{pair.gbm-hand-half.gbm-hand.verdict}) and
\V{pair.gbm-hand-quarter.gbm-hand.gain}
(\V{pair.gbm-hand-quarter.gbm-hand.verdict}). On a quarter of the games GBDT
has the lower log loss (\V{run.gbm-hand-quarter.score}, compared with
\V{run.transformer-medium-hand-quarter.score} for the transformer), and on
half of the games the two models have almost the same log loss
(\V{run.gbm-hand-half.score} and \V{run.transformer-medium-hand-half.score}). The transformer's advantage over
GBDT thus appears only with the full training set (Table~\ref{tab:games},
Figure~\ref{fig:games}).
\fi

\ifmore
\paragraph{Larger training sets.}
The larger sample contains \V{more.gbm-hand.most} the training games, with
the validation and test sets unchanged. Relative to the same model on the
standard training set, \run{transformer-medium-hand} has a paired gain of
\V{pair.transformer-medium-hand-double.transformer-medium-hand.gain} log loss
on twice the training games
(\V{pair.transformer-medium-hand-double.transformer-medium-hand.verdict};
\V{pair.transformer-medium-hand-double.transformer-medium-hand.kept}). For
GBDT the gain is \V{pair.gbm-hand-double.gbm-hand.gain}
(\V{pair.gbm-hand-double.gbm-hand.verdict};
\V{pair.gbm-hand-double.gbm-hand.kept}).

More data helps both model classes. The gain of GBDT is the larger of the
two, but the intervals of the two gains overlap (Table~\ref{tab:games}), so
this experiment does not show which model class benefits more. On twice the
games the transformer remains the better model, with a log loss of
\V{run.transformer-medium-hand-double.score} compared with
\V{run.gbm-hand-double.score} for GBDT, and the intervals of the two models'
gains over the GBDT baseline do not overlap (Table~\ref{tab:leaderboard}).
\ifchangestart
Doubling the training data is worth less to the transformer than pretraining
on the standard training set: the paired gain from pretraining is
\V{pair.transformer-medium-hand-pretrained.transformer-medium-hand.gain}, and
the intervals of the two gains do not overlap (Table~\ref{tab:pairs}).
\fi

The two transformer models take the same number of optimisation steps. On
the standard training set the transformer had its best validation log loss
at epoch \V{run.transformer-medium-hand.best-epoch} of
\V{run.transformer-medium-hand.epochs}. On twice the games it had it in the
final epoch, \V{run.transformer-medium-hand-double.best-epoch} of
\V{run.transformer-medium-hand-double.epochs}, so its validation log loss
was still decreasing when its schedule ended.
\ifchangeepochs
On the standard training set, by contrast, doubling the number of epochs did
not help this model
(\V{pair.transformer-medium-hand-long.transformer-medium-hand.gain}).
\fi
We hypothesise that with more data the transformer is limited by its number
of optimisation steps, in which case its gain from the larger training set
is a lower bound. GBDT was not limited in this way: on twice the games it
used \V{run.gbm-hand-double.trees} of the \V{method.gbm.rounds} boosting
rounds it was allowed.
\fi

% Feature-group ablation: the same ablation for every task, so a model paper
% inputs this file inside its ablation section and then reports its own
% findings. It appears once the task has the runs.
% The prose follows docs/writing.md.
\iffeatures
\subsection{Feature groups}
\label{sec:feature-groups}

\ifsequence
The hand-crafted features fall into three groups: move accuracy, clock use
and the number of available moves.
\else
The hand-crafted features fall into groups: material, mobility, hanging
pieces, passed pawns and king safety, and, where the task gives a line of
moves, the opponent's last move, the first move of the line and the line as a
whole.
\fi
We retrain a model with one group removed at a time
(Table~\ref{tab:features}, Figure~\ref{fig:features}). A significant loss
shows that the model used the group and that no other input provides the same
information. The absence of a loss has two possible explanations: the group
is uninformative, or the same information is available from another group or
from the other inputs. The ablation therefore measures the marginal
contribution of a group given everything else the model receives, which can
be much smaller than the information the group carries on its own.

\begin{table}[htbp]\centering\footnotesize
\block{study-features}
\caption{Feature-group ablation. Each model is trained on all hand-crafted
features except one group. The gain is paired with the model trained on all
features. A negative gain means that the model is worse without the group,
and bold marks statistical significance. Groups are ordered by the size of
the loss.}
\label{tab:features}
\end{table}

\begin{figure}[htbp]\centering
\block{study-features-figure}
\caption{Paired gain of the GBDT baseline without one feature group over the
GBDT baseline with all features, with 95\% CIs. Left of the vertical line,
the model is worse without the group.}
\label{fig:features}
\end{figure}
\fi


\iffeatures
\paragraph{The first move of the solution is the most important group.}
The ablation was run for the GBDT baseline, which is significantly worse in
\V{study.features.gbm-hand.worse.of}. The largest loss follows the removal of
the features that describe the first move of the solution: the piece that
moves, how far it moves, what it captures, whether its destination square is
attacked, and whether the move is a check, a promotion, a retreat or a
sacrifice. Without this group the paired gain is
\V{pair.gbm-hand-no-first-move.gbm-hand.gain}. The first move is the move
that the label is about, so this is the group one would expect to matter
most. The features that describe the solution as a whole are second
(\V{pair.gbm-hand-no-line.gbm-hand.gain}). The length of the solution belongs
to this group, but it is a context feature of this task and is not removed
with it.

Four further groups contribute significantly, and less: the features of the
mover's options, which are whether the mover is in check and the numbers of
legal moves, captures and checks
(\V{pair.gbm-hand-no-mobility.gbm-hand.gain}), the hanging-piece features
(\V{pair.gbm-hand-no-hanging.gbm-hand.gain}), the single feature that states
whether the opponent's mistake was a capture
(\V{pair.gbm-hand-no-setup.gbm-hand.gain}) and the material features
(\V{pair.gbm-hand-no-material.gbm-hand.gain}). Without the features that
count attacks near the kings, GBDT is
\V{pair.gbm-hand-no-kings.gbm-hand.verdict}
(\V{pair.gbm-hand-no-kings.gbm-hand.gain}), and without the passed-pawn
features it is \V{pair.gbm-hand-no-pawns.gbm-hand.verdict}
(\V{pair.gbm-hand-no-pawns.gbm-hand.gain}).

Without the first-move features, GBDT has a log loss of
\V{run.gbm-hand-no-first-move.score}. This is higher than that of the medium
transformer without any hand-crafted features
(\V{run.transformer-medium.score}), which suggests that a network obtains at
least part of the information in this group from the board and the marked
squares of the solution.
\fi
\fi

\section{Error Analysis}
\label{sec:error}

% Error analysis: what the analysis is and how it is computed, the same for
% every task. A model paper inputs this file at the start of its error
% analysis, before the subgroups that are specific to its task.
% The prose follows docs/writing.md.
\iflook
The main results identify the best model. They do not show where a model is
accurate and where it is not, which determines how it can be used. We
therefore analyse the test predictions of every model by subgroup. All
results in this section are computed from the \V{look.test-rows} test
examples (analysis of \V{look.date}).

\paragraph{Subgroups.}
We partition the test set by one variable at a time. For each subgroup we
report the number of test examples, the mean observed target, the mean
prediction of a model and the model's score. Two models are compared:
\V{challenger.name}, the best model other than the GBDT baseline, and the
GBDT baseline itself. The gain of the first over the second is computed
within each subgroup as for the whole test set, as a paired difference over
resampled games (\V{look.resamples} resamples), and is omitted for subgroups
with fewer than \V{look.smallest-compared} test examples. A subgroup whose
interval contains zero provides no evidence in either direction. A numeric
variable is divided into ranges, and a range written $a$-$b$ includes $a$ and
excludes $b$.

\paragraph{Two-way tables.}
For a pair of variables we show the mean observed target and the mean
prediction in each cell, on a common colour scale. Cells with fewer than
\V{grid.smallest-cell} test examples are left empty.

\paragraph{Examples.}
Example \ifsequence games\else positions\fi{} are selected by fixed rules and
not by inspection: typical examples chosen by a hash of their key, the
examples with the best model's most extreme predictions in each direction,
and its largest errors in each direction, \V{examples.count} of each. They
illustrate the quantities under discussion and are not evidence of how often
a pattern occurs.
\else
\emph{This version of the paper was built before the error analysis was run.
The analysis partitions the test set by rating, clock and properties of the
position, measures calibration, and shows example positions. It is produced
by \texttt{lab analyse} on the machine that holds the dataset and the test
predictions of every model, and this section appears once its output is
committed.}
\fi


\iflook
\subsection{Puzzle rating}

The puzzle's rating is not an input of the trained models. The results by
puzzle rating therefore show how much of the difficulty of a tactic the
models infer from the position and the solution. Table~\ref{tab:puzzle-rating}
gives these results, and Figure~\ref{fig:puzzle-rating} compares the observed
share of tactics found in each group with the mean predictions of the two
models.

\begin{table}[htbp]\centering\footnotesize
\block{slice-puzzle-rating}
\caption{Test-set results by the puzzle's rating. For each group: the number
of test positions, the observed share of tactics found, the mean predicted
probability and log loss of each of the two models, and the paired gain of
the first model over the second.}
\label{tab:puzzle-rating}
\end{table}

\begin{figure}[htbp]\centering
\block{slice-puzzle-rating-said}
\caption{Observed share of tactics found and mean predicted probability of
the two models compared, by the puzzle's rating.}
\label{fig:puzzle-rating}
\end{figure}

The share found falls with every group, from
\V{slice.puzzle-rating.-1000.happened} for puzzles rated below 1000 to
\V{slice.puzzle-rating.2200+.happened} for puzzles rated 2200 and above. Both
models follow this trend and underestimate its size. For the easiest group,
the mean predicted probability of \V{challenger.name} minus the observed
share is \V{slice.puzzle-rating.-1000.challenger.off} percentage points, and
for the hardest group it is \V{slice.puzzle-rating.2200+.challenger.off}
points. For the GBDT baseline the differences are
\V{slice.puzzle-rating.-1000.reference.off} and
\V{slice.puzzle-rating.2200+.reference.off}. The models thus overestimate
how often the hardest tactics are found. Among the one-way subgroups of this
section, only those by the time spent on the move show deviations of similar
size (Section~\ref{sec:time}).

\V{challenger.name} is significantly better than the GBDT baseline in
\V{slice.puzzle-rating.ahead.of}, and its gain is largest for the group
\V{slice.puzzle-rating.gain.most.group}
(\V{slice.puzzle-rating.gain.most}). It does not remove the deviation,
however. The puzzle's rating therefore carries information about the
difficulty of a tactic that neither model extracts from the position and the
solution. The application does not have this rating for a tactic in a user's
game, but it could have an estimate of it: a companion task of this
repository predicts the rating of a puzzle from its position and solution.
Whether such an estimate reduces the deviation has not been tested.

\subsection{Properties of the solution}

Several properties of a solution are related to how often it is found. The
shares below depend only on the data and not on any model.
\begin{itemize}
\item \textbf{Sacrifices are missed most often.} We call a first move a
sacrifice when the moved piece can be captured on its destination square, by
a less valuable piece or without being defended, and has captured less than
its own value (the feature \code{first\_sacrifice} of Table~\ref{tab:hand}).
The tactic is found in \V{slice.sacrifice.yes.happened} of the positions with
such a first move, compared with \V{slice.sacrifice.no.happened} of the
others. Such first moves occur in \V{slice.sacrifice.yes.share}\% of the test
positions.
\item \textbf{Checks and captures are found more often.} The share found is
\V{slice.check.yes.happened} when the first move gives check and
\V{slice.check.no.happened} when it does not, and
\V{slice.capture.yes.happened} when the first move is a capture, compared
with \V{slice.capture.no.happened}.
\item \textbf{A retreat makes almost no difference.} The share found is
\V{slice.backward.yes.happened} when the first move goes backward and
\V{slice.backward.no.happened} when it does not.
\item \textbf{The piece that moves first.} Solutions that begin with a king
move are found most often (\V{slice.first-piece.king.happened}) and those
that begin with a pawn move least often
(\V{slice.first-piece.pawn.happened}).
\item \textbf{Longer solutions are found less often.} The share found is
\V{slice.moves.1.happened} for solutions of one move,
\V{slice.moves.2.happened} for two moves, \V{slice.moves.3.happened} for
three and \V{slice.moves.5+.happened} for five or more. The label concerns
only the first move. A plausible explanation of the trend is that the first
move of a long solution is justified only by what follows, so the mover has
to calculate further before playing it.
\item \textbf{Quiet moves are found less often.} The share found is
\V{slice.quiet.0.happened} when every move of the solver is a capture or a
check, and \V{slice.quiet.1.happened} when the solution contains one quiet
move.
\item \textbf{Mates are found more often.} The share found is
\V{slice.mate.yes.happened} when the solution ends in mate and
\V{slice.mate.no.happened} when it does not.
\end{itemize}
The theme tags that Lichess assigns to its puzzles give the same picture
(Table~\ref{tab:slice-themes} in the appendix). The share found is lowest for
the themes \code{sacrifice} (\V{slice.themes.sacrifice.happened}) and
\code{attraction} (\V{slice.themes.attraction.happened}), and highest for
\code{mateIn1} (\V{slice.themes.mateIn1.happened}), \code{oneMove}
(\V{slice.themes.oneMove.happened}), \code{backRankMate}
(\V{slice.themes.backRankMate.happened}) and \code{hangingPiece}
(\V{slice.themes.hangingPiece.happened}).

With respect to these properties and themes both models are much better
calibrated than with respect to the puzzle's rating. Among the groups listed
above, the largest deviation of \V{challenger.name} from the observed share
is \V{slice.first-piece.challenger.widest-off} percentage points, for
solutions that begin with a \V{slice.first-piece.challenger.widest-off.group}
move, and that of the GBDT baseline is \V{slice.themes.reference.widest-off}
points, for the theme \code{\V{slice.themes.reference.widest-off.group}} (the
appendix has a table for each property). The properties, unlike the themes,
are inputs of both models as hand-crafted features, so for them this is
expected.

The gain of \V{challenger.name} over the GBDT baseline is not uniform
(Figure~\ref{fig:themes}). By theme it is significant in
\V{slice.themes.ahead.of} and largest for the theme \code{quietMove}
(\V{slice.themes.quietMove.gain}). By the number of quiet moves in the
solution it rises from \V{slice.quiet.0.gain} with none to
\V{slice.quiet.2.gain} with two. When the first move gives check the gain is
\V{slice.check.yes.gain} (95\% CI \V{slice.check.yes.gain.low} to
\V{slice.check.yes.gain.high}), and when it does not, the gain is
\V{slice.check.no.gain} (\V{slice.check.no.gain.low} to
\V{slice.check.no.gain.high}). For solutions of one move the gain is
\V{slice.moves.1.gain} and not significant. The network's advantage is thus
smallest for short and forcing solutions and largest for solutions with
quiet moves. A plausible explanation is that the hand-crafted features
describe forcing moves almost completely, whereas the value of a quiet move
depends on the rest of the position, which only a network reads.

\begin{figure}[htbp]\centering
\block{slice-themes-gain}
\caption{Paired gain in log loss of \V{challenger.name} over the GBDT
baseline by theme, with 95\% CIs. A puzzle can have several themes.}
\label{fig:themes}
\end{figure}

\subsection{Clock and time spent}
\label{sec:time}

Table~\ref{tab:clock} gives the results by the time remaining on the mover's
clock. With 30 seconds or more, the share of tactics found varies by a few
percentage points: it is \V{slice.clock.30-60.happened} with 30 to 60 seconds,
\V{slice.clock.60-180.happened} with one to three minutes,
\V{slice.clock.180-600.happened} with three to ten minutes and
\V{slice.clock.600+.happened} with ten minutes or more. Below 30 seconds it
falls, to \V{slice.clock.10-30.happened} with 10 to 30 seconds and
\V{slice.clock.-10.happened} with less than ten. In this dataset, time
pressure thus reduces the share of tactics found only in the last half
minute. The two groups below 30 seconds contain \V{slice.clock.-10.share}\%
and \V{slice.clock.10-30.share}\% of the test positions. The comparison is
not controlled for the other differences between these positions, such as the
time control and the stage of the game. The fall below 30 seconds nevertheless appears at every
level of the puzzle's rating (Figure~\ref{fig:grid-clock} in the appendix).

\begin{table}[htbp]\centering\footnotesize
\block{slice-clock}
\caption{Test-set results by the time remaining on the mover's clock, in
seconds.}
\label{tab:clock}
\end{table}

The best model is less well calibrated with respect to the clock than the
GBDT baseline. With less than ten seconds on the clock, the mean predicted
probability of \V{challenger.name} minus the observed share is
\V{slice.clock.-10.challenger.off} percentage points, and with 10 to 30
seconds it is \V{slice.clock.10-30.challenger.off}. For the GBDT baseline the
differences are \V{slice.clock.-10.reference.off} and
\V{slice.clock.10-30.reference.off}. \V{challenger.name} is significantly
better than the GBDT baseline in \V{slice.clock.ahead.of}
(\V{slice.clock.ahead.groups}), and with less than ten seconds its gain is
\V{slice.clock.-10.gain} and not significant. The results by game speed show
the same pattern (Table~\ref{tab:slice-speed} in the appendix). The share
found is \V{slice.speed.blitz.happened} in blitz games,
\V{slice.speed.rapid.happened} in rapid games and
\V{slice.speed.classical.happened} in classical games, and the deviations of
\V{challenger.name} are \V{slice.speed.blitz.challenger.off} points for blitz
and \V{slice.speed.classical.challenger.off} points for classical games. The
network's predictions thus depend less on the time available than the
observed shares do. We have not identified the cause.

The time spent on the move is known only after the move and is not an input.
Table~\ref{tab:seconds} and Figure~\ref{fig:seconds} show that it is strongly
related to the outcome. The share found is \V{slice.seconds.-2.happened} for
moves played in less than two seconds and \V{slice.seconds.2-5.happened} for moves
that took two to five seconds. It then falls with every group, to
\V{slice.seconds.40+.happened} for moves that took 40 seconds or more.
Tactics are thus found more often on moves that are played quickly, and after
a long think the mover has more often chosen another move.

\begin{table}[htbp]\centering\footnotesize
\block{slice-seconds}
\caption{Test-set results by the time spent on the move, in seconds. The
time spent is known only after the move and is not an input of any model.}
\label{tab:seconds}
\end{table}

\begin{figure}[htbp]\centering
\block{slice-seconds-said}
\caption{Observed share of tactics found and mean predicted probability of
the two models compared, by the time spent on the move.}
\label{fig:seconds}
\end{figure}

Both models reproduce this relation only in part. For moves that took 40
seconds or more, the mean predicted probability of \V{challenger.name} minus
the observed share is \V{slice.seconds.40+.challenger.off} percentage points,
and for moves that took two to five seconds it is
\V{slice.seconds.2-5.challenger.off}. For the GBDT baseline the differences
are \V{slice.seconds.40+.reference.off} and
\V{slice.seconds.2-5.reference.off}. The time spent on a move is thus
informative about the outcome beyond what either model predicts from its
inputs. A companion task of this repository predicts
this time from a position, and a model trained on both targets is a natural
experiment.

\subsection{Rating of the mover and of the puzzle}

By band of the mover's rating, the share of tactics found rises from
\V{slice.band.-1700.happened} for movers rated below 1700 to
\V{slice.band.2100-2299.happened} for movers rated 2100 to 2299. The mean
predictions of \V{challenger.name} and of the GBDT baseline are within
\V{slice.band.challenger.widest-off} and \V{slice.band.reference.widest-off}
percentage points of the observed share in every band
(Table~\ref{tab:slice-band} in the appendix). The difference between the
bands is small beside the difference between easy and hard puzzles.
Figure~\ref{fig:grid-ratings} shows both variables together.

\begin{figure}[htbp]\centering
\block{grid-ratings}

\medskip
\block{grid-ratings-said}
\caption{Share of tactics found in percent by the mover's rating and the
puzzle's rating. Top: observed shares. Bottom: mean predictions of the best
model other than the GBDT baseline. Both panels use the same colour scale.}
\label{fig:grid-ratings}
\end{figure}

Within every band of the mover's rating, the share found falls steeply from
the easiest to the hardest puzzles, for example from
\V{grid.ratings.-1700.-1000.happened} to
\V{grid.ratings.-1700.2200+.happened} for movers rated below 1700 and from
\V{grid.ratings.2300+.-1000.happened} to
\V{grid.ratings.2300+.2200+.happened} for movers rated 2300 and above. Within
a column, it rises with the mover's rating by much less. For puzzles rated
1800 to 2200 it rises from \V{grid.ratings.-1700.1800-2200.happened} to
\V{grid.ratings.2300+.1800-2200.happened}, and for puzzles rated below 1000
it is no higher for the strongest movers than for the weakest. The share
found is therefore not a function of the difference between the two ratings
alone, as the Elo formula assumes, which is one reason for the poor fit of
the rating-based baseline (Section~\ref{sec:rating-baseline}).

The model reproduces the pattern and compresses it along the puzzle's
rating. Its mean prediction is above the observed share in every cell of the
two hardest columns, and below it in every cell of the two columns from 1000
to 1800. For the hardest puzzles the observed share hardly rises with the
mover's rating, from \V{grid.ratings.-1700.2200+.happened} to
\V{grid.ratings.2300+.2200+.happened}, whereas the prediction rises from
\V{grid.ratings.-1700.2200+.said} to \V{grid.ratings.2300+.2200+.said}. The
largest deviation in a cell is
\V{grid.ratings.transformer-medium-hand-pretrained.widest-off} percentage
points, for movers rated
\V{grid.ratings.transformer-medium-hand-pretrained.widest-off.row} and
puzzles rated
\V{grid.ratings.transformer-medium-hand-pretrained.widest-off.column},
compared with \V{grid.ratings.gbm-hand.widest-off} for the GBDT baseline. The
rating-based baseline, which receives the puzzle's rating, deviates by up to
\V{grid.ratings.puzzle-rating.widest-off} points, in the same cell.

\subsection{Win probability before the mistake}

The mover's win probability before the opponent's mistake is an input. The
mean predictions of \V{challenger.name} and of the GBDT baseline are within
\V{slice.prior.challenger.widest-off} and
\V{slice.prior.reference.widest-off} percentage points of the observed share
in every group (Table~\ref{tab:slice-prior} in the appendix). The share of tactics found is
highest when the mover was already winning before the opponent's move:
\V{slice.prior.90+.happened} at a win probability of 90\% or more, a group
that contains \V{slice.prior.90+.share}\% of the test positions. It is lowest
at win probabilities between 30\% and 50\% (\V{slice.prior.30-50.happened})
and between 10\% and 30\% (\V{slice.prior.10-30.happened}).
\V{challenger.name} is significantly better than the GBDT baseline in
\V{slice.prior.ahead.of} (\V{slice.prior.ahead.groups}).

\subsection{Agreement with the GBDT baseline}

Two models with similar scores can still disagree on individual positions, in
which case an average of their predictions may be more accurate than either.
Table~\ref{tab:beside} gives, for each model, the correlation of its
predicted probabilities with those of the GBDT baseline. For the best model
the correlation is \V{best.beside.correlation}. The two models therefore
disagree on many individual positions, which suggests that an ensemble of the
two could improve on both. We have not trained one. The table also gives the
share of test positions on which a model's prediction is closer to the
outcome than the baseline's, which is \V{best.beside.closer}\% for the best
model.

\begin{table}[htbp]\centering\small
\block{beside}
\caption{Agreement with the GBDT baseline on the test set: the correlation
between the two models' predicted probabilities, and the share of positions
on which the model's prediction is closer to the outcome than the
baseline's.}
\label{tab:beside}
\end{table}

\subsection{Example positions}
\label{sec:positions}

The following examples were selected from the test set by the fixed rules
described above. Each board is shown from the mover's side. The grey arrow is
the opponent's preceding move, the blue arrow is the first move of the
solution, and an orange arrow is the move that was played where it differs
from the solution.
Next to each board we give the opponent's move, the solution, the move that
was played, whether the tactic was found, the probabilities that the best
model and the GBDT baseline predicted, and the context of the position. The
context includes the puzzle's rating, its themes and the time spent on the
move, which the models of the fixed sequence do not receive. The positions
are Lichess puzzles and are identified by their puzzle identifiers.

Figure~\ref{fig:typical} shows three typical positions. The tactic was found
in all three. The predictions of the two models lie between 62\% and 81\%,
and in each position they differ from each other by ten percentage points or
more, which illustrates the disagreement between the two models described
above.

The largest errors of a probabilistic model are of two kinds:
Figure~\ref{fig:surprised-no} shows the positions with the highest predicted
probability in which the tactic was missed, and
Figure~\ref{fig:surprised-yes} shows those with the lowest predicted
probability in which it was found. Neither kind is necessarily a mistake of
the model. A predicted probability of 99\% is correct if one such tactic in a
hundred is missed, and these examples are the misses.

In the three positions of Figure~\ref{fig:surprised-no} the solution is a
mate in one or two moves, and both models predicted a probability above
98\%. In each of them the mover played another move in four seconds or less.
A model of the position cannot foresee such a lapse, and predictions of this
size are not wrong on account of it: among the 1\% of test positions with the
highest predictions, \V{best.top.1.rate}\% of the tactics were found. In the
three positions of Figure~\ref{fig:surprised-yes} the solution has three
moves and begins with a sacrifice, in which the piece that moves first is
captured at once, and the puzzles are rated between 1996 and 2604. The best
model predicted probabilities between 5\% and 7\%, less than half of those
of the GBDT baseline, and the tactic was found each time. The positions with
the highest and the lowest predictions overall are shown in the appendix
(Figures~\ref{fig:highest} and~\ref{fig:lowest}).

\examples{examples-typical}{Typical positions: the first \V{examples.count}
test positions in the order of a hash of their key.}{fig:typical}

\examples{examples-surprised-no}{High predicted probability and a missed
tactic: the test positions with the highest predictions of the best model
among those in which the tactic was missed.}{fig:surprised-no}

\examples{examples-surprised-yes}{Low predicted probability and a tactic
that was found: the test positions with the lowest predictions of the best
model among those in which the tactic was found.}{fig:surprised-yes}
\fi

\section{Practical Use}
\label{sec:use}

\begin{itemize}
\item \textbf{The prediction can be shown as a probability.} After
calibration, the mean predicted probability of the best model is within
\V{best.calibration.widest-gap} percentage points of the observed share of
tactics found in every bin, so for Lichess games like those of the dataset
the number means what it says on average.\iflook{} The exception is the
hardest tactics: for puzzles rated 2200 and above the mean prediction is
\V{slice.puzzle-rating.challenger.widest-off} points above the observed
share.\fi
\iflook
\item \textbf{Distinguishing a lapse from a tactic beyond the player's
level.} Among the 10\% of test positions with the highest predictions,
\V{best.top.10.rate}\% of the tactics were found, so a miss there is a
lapse. Among the 10\% with the lowest predictions, \V{best.bottom.10.rate}\%
were missed, so a miss there is what usually happens at that level.
\item \textbf{Choosing what to practise.} The tactics that players miss most
often in their games begin with a sacrifice
(\V{slice.sacrifice.yes.happened} found), contain quiet moves or have long
solutions. Mates in one, back-rank mates and hanging pieces are among the
tactics found most often (\V{slice.themes.mateIn1.happened},
\V{slice.themes.backRankMate.happened} and
\V{slice.themes.hangingPiece.happened}).
\item \textbf{Time pressure matters late.} The share of tactics found
changes little down to 30 seconds on the clock and falls below that. The
best model overestimates the share for these positions: with less than ten
seconds its mean prediction is \V{slice.clock.challenger.widest-off} points
above the observed share.
\item \textbf{The puzzle's rating is not a substitute.} The Elo formula
applied to the mover's rating and the puzzle's rating has a log loss of
\V{run.puzzle-rating.score}, compared with \V{reference.score} for the GBDT
baseline, which does not use the puzzle's rating.
\fi
\item \textbf{GBDT provides most of the accuracy.} The GBDT baseline has a
log loss of \V{reference.score}, compared with \V{best.score} for the best
model and \V{floor.score} for the constant baseline, and needs no GPU. Where
the application has to run the model itself, GBDT is the natural starting
point.
\end{itemize}

\section{Limitations}
\label{sec:limits}

\begin{itemize}
\item \textbf{The solution is an input.} The models are told the tactic. In
the application the solution would come from an engine, and a line that an
engine finds in an arbitrary game need not resemble the solutions of
published puzzles.
\item \textbf{Only the first move is labelled.} A mover who plays the first
move without having seen the rest is counted as having found the tactic, and
a mover who plays another move that keeps the advantage is counted as having
missed it.
\item \textbf{Only tactics that work.} Every position of the dataset is a
published puzzle, so its tactic works. The models say nothing about positions in which a tempting
sacrifice fails, and such false alarms are part of what a player has to
judge in a game.
\item \textbf{The positions are a selected sample.} They are positions that
Lichess published as puzzles and whose puzzles have an established rating.
They are not a random sample of the tactics that occur in games.
\iflook
\item \textbf{Few weaker players.} Only 5\% of the movers are rated below
\V{describe.mover-elo.q5}, so the results say little about weaker players.
\item \textbf{The hardest tactics are overestimated.} For puzzles rated 2200
and above, the mean prediction is
\V{slice.puzzle-rating.challenger.widest-off} percentage points above the
observed share.
\fi
\item \textbf{No bullet games.} Bullet games are excluded, so the models say
nothing about the fastest games.
\item \textbf{A single position without history.} The models do not see the
moves before the opponent's mistake or how the players used their clocks
before the position.
\item \textbf{No comparison with a move-prediction model.} The probability
that a model of human moves~\cite{maia2} assigns to the first move of the
solution is a natural alternative that we have not evaluated.
\item \textbf{Single training runs.} Each network was trained with one
random seed. The confidence intervals reflect the sampling of test games and
not the variation between training runs.
\end{itemize}

\section{Reproducibility}
\label{sec:cost}

Training the \V{runs.count} models on the leaderboard took \V{cost.seconds}
in total: \V{cost.seconds.networks} for the networks, each trained on a
single cloud GPU, and \V{cost.seconds.others} for the baselines and GBDT on a
CPU. The ablation models that are not on the leaderboard are not included in
this total. The best model has \V{best.parameters} parameters and was trained
for \V{best.epochs.run} epochs in \V{best.seconds} on an \V{best.device}, at
about \V{best.rows-a-second} positions per second, not counting the
pretraining of its encoder. Table~\ref{tab:runs} lists the hardware and
training time of every model, and Table~\ref{tab:settings} the
hyperparameters of the networks of the fixed sequence.

All results are generated from committed files. The dataset is built from
its recipe by \texttt{lab data tactics}, and the recipe names the revision of
the source that a build reads. The models are trained by
\texttt{lab all tactic-found}, each training run writes a record to the
model's \texttt{results/} directory, and \texttt{lab paper tactic-found}
generates every number, table and figure of this paper from those records
and compiles it. A record names the build of the dataset it used
(\V{data.fingerprint}) and the commit of the code that produced it
(\V{code.commit}). Model weights and test predictions are not stored in the
repository. The model's card states where they are kept and gives the hash
of the best model's weights.

\section{Conclusion and Future Work}
\label{sec:conclusion}

Whether a player finds a tactic in a game can be predicted with a log loss of
\V{best.score}, compared with \V{floor.score} for a constant prediction. The
largest part of this improvement comes from hand-crafted features of the
position and of the solution, and a smaller part from the ratings, the clocks
and the evaluations. A network that reads the board improves on GBDT
only when it is given the hand-crafted features or a pretrained encoder. The
predicted probabilities are calibrated over the whole test set. They are too
high for the hardest tactics, for moves that took long and for positions
with little time on the clock. The following experiments are the natural
next steps.
\begin{itemize}
\item \textbf{An estimate of puzzle difficulty as an input.} The models are
miscalibrated with respect to the puzzle's rating, which the application
does not have. A predicted rating from the companion task, or an encoder
pretrained on that task, may supply part of the missing information.
\ifchangestart
\item \textbf{A stronger pretrained encoder.} The encoder was pretrained
once, at one size, on part of the evaluation database. Longer pretraining,
more positions and a larger encoder have not been tried.\ifchangeconv{} A
pretrained ResNet has not been tried either, although the ResNet is the best
model trained from random initialisation.\fi
\fi
\ifmore
\item \textbf{More data together with more optimisation steps.} On twice the
games the transformer's validation log loss was still decreasing when its
schedule ended. Training on the larger set with a longer schedule, and
training the pretrained model on it, are the two missing cases.
\fi
\item \textbf{A label for the whole solution.} The source games show whether
the mover went on to play the rest of the solution, and their evaluations
show whether a different move kept the advantage. Both would make the label
closer to what a coach means by finding a tactic.
\item \textbf{A joint model of tactics and thinking time.} The time spent on
a move is related to whether the tactic is found in a way that the present
models do not capture.
\item \textbf{Weaker players.} A dataset of tactics from the games of weaker
players would cover the users whom the present data underrepresent.
\item \textbf{False alarms.} Positions in which a tempting tactic fails
would allow a model to predict errors of commission as well as omissions.
\item \textbf{A baseline from a move-prediction model.} A human-move model
gives a probability for the first move of the solution without being told
that it is one.
\item \textbf{An ensemble.} The best network and the GBDT baseline disagree
often enough that an average of the two may be more accurate than either, at
the cost of running both.
\end{itemize}

\bibliographystyle{unsrt}
\bibliography{../../../papers/references}

\clearpage
\appendix

\section{Hyperparameters and Features}

\begin{table}[htbp]\centering\footnotesize
\block{settings}
\caption{Hyperparameters of the networks of the fixed sequence, from their
configuration files. \emph{LR} is the peak learning rate. Values shared by
all of them are given below the table.}
\label{tab:settings}
\end{table}

\begin{table}[htbp]\centering\footnotesize
\block{hand}
\caption{Hand-crafted features: the inputs of \V{task.reference} in addition
to the scalar context features, and of a network whose configuration includes
them.}
\label{tab:hand}
\end{table}

\section{Calibration by Bin}

\begin{table}[htbp]\centering\small
\block{calibration-table}
\caption{Calibration by bin on the test set, for the best model and the GBDT
baseline. The bins of each model are formed from its own predicted
probabilities and have equal size. \emph{Predicted} is the mean predicted
probability in a bin, and \emph{observed} is the share of tactics in the bin
that were found.}
\label{tab:calibration}
\end{table}

\iflook
\section{Clock and Puzzle Rating}

\begin{figure}[htbp]\centering
\block{grid-clock-rating}

\medskip
\block{grid-clock-rating-said}
\caption{Share of tactics found in percent by the time remaining on the
mover's clock, in seconds, and the puzzle's rating. Top: observed shares.
Bottom: mean predictions of the best model other than the GBDT baseline. Both
panels use the same colour scale.}
\label{fig:grid-clock}
\end{figure}

\section{Examples with Extreme Predictions}

The three positions with the highest predictions are back-rank mates in one
move whose puzzles are rated between 600 and 622, and all three were found
(Figure~\ref{fig:highest}). The three positions with the lowest predictions
have solutions of three moves and puzzle ratings above 2000. Two of them
begin with a queen sacrifice and the third is a pawn endgame, and all three
were missed (Figure~\ref{fig:lowest}). The best model predicted
probabilities between 3\% and 4\% for these three positions, and the GBDT
baseline between 8.6\% and 11.2\%.

\examples{examples-highest}{The test positions with the highest predicted
probabilities of the best model.}{fig:highest}

\examples{examples-lowest}{The test positions with the lowest predicted
probabilities of the best model.}{fig:lowest}

\section{Subgroup Tables}

Each table below partitions the test positions by one variable. The columns
are the group, the number of test positions, the observed share of tactics
found, the mean predicted probability and log loss of each of the two models
compared, and the paired gain of the first model over the second with its
95\% CI, in bold when statistically significant. A puzzle can have several
themes, so the groups of the table by theme overlap.

\block{slices}
\fi

\end{document}
